\documentclass[11pt,twoside,openany]{book}
\usepackage[paperwidth=450pt,paperheight=666pt,inner=48pt,outer=38pt,top=40pt,bottom=44pt,headheight=14pt,headsep=13pt]{geometry}
\usepackage{amsmath,amssymb,graphicx,booktabs,longtable,array,xcolor,hyperref,fancyhdr,titlesec,microtype}
\definecolor{ink}{HTML}{173E48}
\definecolor{quiet}{HTML}{56666A}
\hypersetup{colorlinks=true,linkcolor=ink,citecolor=ink,urlcolor=ink,pdftitle={Why EBU? Life, Balance, and the Accounting of Physical Action},pdfauthor={Konrad Grzyb},pdfsubject={New-source Book I introductory edition; candidate for review}}
\urlstyle{same}
\setlength{\parindent}{12pt}
\setlength{\parskip}{2pt}
\linespread{1.04}
\setcounter{tocdepth}{0}
\titleformat{\chapter}[display]{\color{ink}\normalfont}{\small\sffamily BOOK I\quad /\quad CHAPTER \thechapter}{8pt}{\raggedright\hyphenpenalty=10000\fontsize{21}{24}\selectfont}
\titlespacing*{\chapter}{0pt}{0pt}{18pt}
\titleformat{\section}{\color{ink}\normalfont\large}{\thesection}{0.65em}{}
\titlespacing*{\section}{0pt}{14pt plus 3pt minus 2pt}{6pt}
\titleformat{\subsection}{\normalfont\bfseries}{\thesubsection}{0.6em}{}
\pagestyle{fancy}\fancyhf{}
\fancyhead[LE,RO]{\small\sffamily\thepage}
\fancyhead[LO]{\small\sffamily Why EBU?}
\fancyhead[RE]{\small\sffamily\nouppercase{\leftmark}}
\renewcommand{\chaptermark}[1]{\markboth{Chapter \thechapter}{}}
\renewcommand{\headrulewidth}{0.3pt}
\fancypagestyle{plain}{\fancyhf{}\fancyfoot[C]{\small\sffamily\thepage}\renewcommand{\headrulewidth}{0pt}}
\newcommand{\question}[1]{\noindent\textit{#1}\par\medskip}
\newcommand{\carry}[1]{\par\medskip\noindent\textbf{Carry forward.} #1\par}
\newcommand{\practice}[2]{\par\medskip\noindent\textbf{Pause and try.} #1\par\noindent\textit{A worked response.} #2\par}
\newcommand{\illustration}[3]{\begin{figure}[!htbp]\centering\includegraphics[width=\linewidth]{figures/#1.pdf}\caption{\small #2}\label{#3}\end{figure}}
\newcommand{\unit}{\mathrm{L}}
\setlength{\emergencystretch}{1.4em}
\widowpenalty=10000\clubpenalty=10000
\raggedbottom
\renewcommand{\topfraction}{0.9}
\renewcommand{\textfraction}{0.08}
\renewcommand{\floatpagefraction}{0.85}
\setlength{\intextsep}{10pt}
\let\cleardoublepage\clearpage
\makeatletter
\renewcommand{\l@chapter}[2]{\@dottedtocline{0}{0pt}{1.8em}{#1}{#2}}
\makeatother
\begin{document}
\frontmatter
\begin{titlepage}
\color{ink}
\noindent\small\sffamily ENERGY BALANCE PROJECT\par
\vspace{49pt}
\noindent\fontsize{42}{45}\selectfont Why EBU?\par
\vspace{20pt}
\noindent\fontsize{23}{29}\selectfont Life, Balance, and the\par\noindent Accounting of Physical Action\par
\vspace{32pt}\hrule\vspace{21pt}
\noindent\large Book I\par
\vspace{8pt}\noindent\normalsize A new-source introductory edition\par
\vfill
\noindent\large Konrad Grzyb\par
\vspace{20pt}
\noindent\small\sffamily CANDIDATE FOR AUTHOR REVIEW\par
\vspace{5pt}\noindent 19 September 2026\par
\end{titlepage}
\chapter*{About this edition}
Konrad Grzyb is the author and originator of the Energy Balance Project. This candidate was researched, drafted and typeset with AI assistance under the author's instruction. It is presented for review; its delivery does not imply the author's acceptance, external peer review or publication.

This is a new manuscript. The preserved first volume, \textit{Physical Intuition and the EBU Idea}, provides a content baseline, but its original editable sources and complete build system have not been recovered. The new teaching order follows the project's September 2026 introductory blueprint. Original Parts I, II and III remain unchanged. Chapter numbers in this edition are new: a reference to historical Part II.33 still means Chapter 33 of the preserved Part II.

The software repository has a separate MIT licence. That software licence is not a licence for reproducing the preserved books. This candidate introduces no new rights grant. Source identities, content dispositions and build instructions accompany the editable manuscript.

The book separates observed evidence, definitions, conditional mathematics, illustrative examples, proposed rules and open questions. It reports no new scientific experiment. Numerical tank examples are arithmetic teaching devices, not field measurements or a registered experimental parameter set. Every illustration is labelled by its status.

The Local Gaussian model is a prospective research direction. Its equations can be explained and checked before a new experiment. Its practical usefulness, recovery behaviour and institutional consequences cannot be established by explanation alone.
\chapter*{An invitation to the reader}
A clinic can pay its bills and still face a broken pump. It can also have a working pump while the people who need its water cannot reach it. These are different failures. To understand them, we need to keep physical function, institutional claims and human access in view at the same time.

This book asks whether a precise account of physical actions could contribute something useful to that task. It does not begin by asking the reader to accept a new economy. It begins by asking what a record would have to say before it deserved our trust.

The first five chapters develop the reason for the question. Chapters 6--9 introduce functioning systems, physical boundaries and actions. Chapters 10--13 build a declared ruler, one step at a time. Chapters 14--17 use it to calculate finite changes and explain exactly which identities follow. Chapter 18 returns to the clinic and the work still ahead.

No calculus, graph theory or programming is assumed. When a new symbol appears, first read its name and unit. Follow the arithmetic before trying to remember the formula. A superscript may mean a square, a reference label or a transpose; the text will say which. The glossary and symbol table are there for returning to an idea, not for memorising them before beginning.

The recurring tank example is deliberately small. Its simplicity lets us see where a sign comes from, where a term cancels and where an assumption enters. We will not pretend that two water stocks represent an entire clinic, much less the planet.

A useful outcome is the ability to ask a better question of any EBU claim: what changed, relative to which reference, inside which boundary, and with what evidence?
\tableofcontents
\mainmatter
\chapter{The question beneath an economy}\label{ch:question}
\question{What are economic systems ultimately meant to make possible?}

The water arrives before the first patient. Someone opened a valve, a motor turned, and a pump lifted water into the clinic's tank. A nurse can wash her hands. Instruments can be cleaned. The building is able to do work that people depend on.

Several records describe the morning. An invoice records electricity purchased. A wage record identifies payment to the person who maintained the pump. A contract assigns responsibility for repairs. A meter records water delivered. An inspection records the condition of a seal. None of these records is foolish or dispensable. Each answers a different question.

Now suppose the electricity invoice is paid on time, but the seal leaks. The payment record may be perfectly correct while the physical service deteriorates. Alternatively, the pump may work perfectly while a patient lacks transport to the clinic. The engineering record can be correct while access fails. Good accounting begins by refusing to turn one kind of success into every kind of success.

\section{Claims, things and services}
An economic claim gives someone a recognised entitlement within an institutional arrangement: to payment, repayment, delivery or use. A physical resource is something represented in the material world, such as water in a tank. A service is what an arrangement makes possible, such as safe washing at a particular time. A human need explains why that service matters.

These descriptions connect, but do not replace one another. A contract for water can help organise a reliable supply. It cannot itself hydrate a patient. A tank of water may help provide care, but water of the wrong quality, in the wrong place, or available after the emergency is not the same service. A litre is not a promise, and a promise is not a litre.

The distinction is especially important for maintenance. A replacement bearing can be physically small and functionally decisive. The account must say what it represents before assigning a number. Does it record the bearing's mass, the motor's operating condition, the electricity consumed during repair, or the increased ability to deliver water later? Those are different objects, with different evidence requirements.

\illustration{two_records}{Schematic: one pumping event supports several records. The payment and physical records can be linked by event identity without being interchangeable. No EBU valuation is assigned by this diagram.}{fig:records}

\section{What EBU proposes to ask}
The Energy Balance Project investigates an explicit account of physical action. In this book, an Energy Balance Unit, or EBU, expresses a signed change in a declared measure of represented burden, after any separately justified process burden. The declaration matters: an account has to identify what it measures and what it leaves outside.

Before we can calculate, we must answer ordinary questions. What was present before the action? What did the action change? Which physical restrictions applied? Relative to what reference is the change evaluated? Has an important effect already been counted elsewhere? Who chose the action? Who verified that it occurred?

The resulting number cannot contain all possible answers about the clinic. It need not be a price. It does not state the moral worth of the nurse, mechanic or patient. It cannot certify an effect that the chosen state description does not represent. A local equation is useful only if its limited meaning stays visible.

\section{The ethical starting point}
The motivation of this book includes an ethical commitment: people should have reliable access to essential support, and avoidable damage to the conditions of life deserves attention. That is a declared value commitment. It is not a fact proved by a mathematical curve.

Different communities may disagree about priorities, responsibilities and acceptable trade-offs. A transparent physical account could inform those disagreements. It cannot eliminate them by renaming a disputed choice a measurement. The right to receive care must not be silently made conditional on a person's ability to generate a positive score in a model.

The central research argument is therefore precise. Monetary transactions and existing institutions do not automatically guarantee preservation of all specified life-supporting functions. EBU investigates an explicit physical-action account; its usefulness remains a research question. Evidence that a problem exists is a reason to investigate. It is not evidence that this proposal solves it.

\practice{The clinic paid for ten units of water, but its meter recorded eight. Which record settles the physical delivery question?}{The meter is relevant to the delivered quantity, subject to its accuracy and placement. The invoice establishes what was charged or contracted. Comparing the records may reveal a discrepancy, but neither alone tells us whether two units leaked, were stored elsewhere, or were mismeasured. A missing explanation should remain visible.}

\carry{We need a physical account because physical functioning is one of the things we care about. To understand its proposed role, we should first understand the considerable work monetary signals already do.}

\chapter{What monetary signals can and cannot tell us}\label{ch:money}
\question{Why can a financially attractive action have harmful consequences outside its price?}

The clinic's invoice is not an illusion. Payment lets a supplier meet obligations, purchase replacement parts and pay workers. Prices allow many unlike goods to be compared within a common monetary account. A household can make a budget; a firm can examine revenue and expense; a public provider can plan expenditure. These are practical achievements.

A price is also a signal formed within institutions. It can change when production becomes harder, demand rises, alternatives disappear, contracts change or market power shifts. Its meaning depends on who can offer, who can buy and what rules apply. A signal produced by these conditions need not be a complete description of the physical consequences of an action.

\section{Purchasing power and credit}
Purchasing power concerns what a person or institution can obtain with available monetary resources at prevailing prices. It is different from need. Two people may need the same essential service while having very different capacities to pay. A price records the terms of an exchange; it does not directly measure the urgency of every excluded person's need.

Credit connects present spending to future obligations. In its explanation of modern UK banking, the Bank of England describes how a commercial-bank loan creates a matching deposit. It also explains constraints on lending, including profitability, regulation and monetary policy. This is an institutional mechanism, not unlimited creation of physical wealth.\cite{boe}

Imagine a loan financing the clinic's pump. It can make an order possible today, ahead of future revenues. The pump still requires manufacture, transport, installation and maintenance. Creating a claim and supplying the physical capacity are connected activities, but only the second installs equipment. Conversely, refusing a loan can leave usable technical possibilities unrealised. The physical and financial questions remain linked without becoming identical.

\section{Consequences outside the transaction}
An externality is an effect on people outside a transaction that is not fully reflected in the decision-makers' private costs or benefits. Public goods pose a related problem when people cannot readily be excluded from the benefit and one person's use does not subtract from another's. These concepts explain why private incentives can differ from wider consequences; they do not say that every market always fails.\cite{externalities}

Consider a hypothetical workshop choosing whether to contain a harmful discharge. If its invoice includes the cost of disposal but excludes damage downstream, the cheapest option for the workshop can burden others. The example assumes an unpriced effect. It does not establish how common that situation is or measure its total cost.

Change the rules and the decision can change. A discharge limit, an enforceable liability rule or a charge tied to emissions can make prevention part of the workshop's own calculation. Public provision can organise a service that individual purchases would undersupply. Rules need information and enforcement, but their existence matters to the argument.

The EU Emissions Trading System supplies an actual example of a connection between a physical quantity and a monetary incentive: covered operators monitor emissions and surrender allowances within a cap-and-trade system. This does not establish that the instrument covers every ecological effect or that its price is adequate for every purpose. It does refute the claim that environmental effects never enter monetary institutions.\cite{ets}

\section{A price and a physical history}
Suppose two deliveries cost the clinic the same amount. One uses a well-maintained pipe, the other an emergency tanker. Equal payments do not logically imply equal fuel use, water loss, reliability or exposure to disruption. To answer those questions, we need measurements and a boundary for the comparison.

The reverse also holds. Two physically similar deliveries may have different prices because ownership, taxes, contracts or bargaining conditions differ. A physical account does not determine those legal terms by itself. Neither price nor physical description has to be declared meaningless for their difference to matter.

There are deliberate feedback institutions already. The ECB's 2025 strategy names price stability as its primary objective and commits to considering climate and nature implications within its mandate. A stated monetary mandate is not a guarantee of every social or ecological function, but it is plainly a feedback arrangement.\cite{ecb}

For EBU, the question is narrower than whether money is good or bad. Could a declared physical-action account make certain consequences explicit in a way that is useful, auditable and sufficiently local? It would have to earn its place alongside existing accounts and rules. Calling a number physical does not automatically make it complete or well calibrated.

\practice{A pollution charge makes the workshop's cleaner option cheaper. Does this make a physical record redundant?}{No. The charge itself depends on a defined effect and some means of measurement or estimation. The record and incentive can support each other. Nor does the charge's existence prove all downstream effects have been included.}

\carry{Omitted effects and inability to purchase essentials are different problems. The next chapter asks why a resource can exist while a person's access to it remains inadequate.}

\chapter{Scarcity, access and unequal exposure}\label{ch:access}
\question{Why can needs remain unmet even when some relevant resources exist?}

A dry well, a broken pipe and an unpaid connection fee can all prevent water from reaching a home. They are not the same obstacle. More money cannot instantly refill the well. More rainfall cannot by itself repair the pipe. A working pipe does not settle who is allowed to use it or who can afford the service.

FAO distinguishes physical scarcity, access scarcity and infrastructure scarcity in its account of water. The distinction keeps the availability of suitable water separate from the arrangements needed to deliver it.\cite{water} We can extend the question cautiously in our clinic example: what exactly prevents a specified service from reaching a specified person at a specified time?

\section{Finding the constraint}
Physical availability includes quantity, location, timing and quality. A region may have water in one season and inadequate storage for another. Water can be present but unsuitable for a particular use. In our example these are questions for physical measurement and domain knowledge, not judgments inferred from the patient's bank balance.

Infrastructure includes treatment, pipes, pumps, roads, power and the ability to maintain them. An institution can own abundant equipment that it cannot operate because a small component is missing. Distribution concerns where usable resources go and under what rules. Affordability concerns whether the required payment fits within a person's resources and competing essential expenses.

These constraints may reinforce one another. A broken network may force people to purchase a more expensive substitute. An unaffordable service may remain unused while installed capacity stands idle. Neither observation would prove that physical limits are unreal. It would tell us to identify the active constraint before choosing a remedy.

Food provides the same conceptual distinction without requiring a disputed global count. Food on a warehouse shelf is physical availability at that place. A household's ability to obtain, store, prepare and safely consume it involves further conditions. A delivery record and a purchase record each answer part of that story. Neither is a complete measure of nourishment.

\section{Socioeconomic conditions are several things}
Socioeconomic status, often shortened to SES, describes a person's or household's position through measures such as income, education and occupation. Wealth, housing security, social exclusion and health add related but distinct information. An income number alone does not describe all of these conditions.

The WHO's 2025 report on social determinants of health equity describes how living and working conditions and access to power and resources influence avoidable health gaps.\cite{health} That supports attention to mechanisms. It does not establish that every health difference has a single monetary cause.

Cause and association must also be distinguished. Poor health can reduce earning opportunities; insecure work can make it harder to maintain health; earlier conditions can influence both. A comparison between groups does not automatically tell us what caused one individual's circumstances. To identify an effect, a study needs a design that can address competing explanations.

One bounded example is the US Moving to Opportunity experiment. Families in five cities were randomly offered different housing-assistance options during 1994--1998. A later study found improved adult outcomes for children who moved young, while effects for adolescents differed. This is evidence about a particular intervention and population, not a universal coefficient linking money to opportunity or a test of EBU.\cite{mto}

\section{Two households in the same heatwave}
Imagine two households during the same hot week. One has shade, reliable electricity, a room that stays cooler and flexible work. The other lives under an uninsulated roof, has a member who needs assistance and cannot easily leave work or travel. This is a schematic example of unequal exposure and ability to respond, not a clinical risk calculation.

The same outdoor temperature does not make their situations identical. A physical record limited to outdoor temperature would omit the housing and access differences. A record limited to income would omit other relevant conditions. Expanding a model can make it more informative, but also creates demands for evidence, privacy protection and clear interpretation.

The example illustrates a difficulty for any scalar account. If a person's disability makes a service physically more demanding to provide, a higher recorded burden must not become a verdict that the person deserves less support. Recording resource use and deciding entitlements are separate tasks. The commitment to essential access needs an institutional expression of its own.

\section{What an account could contribute}
A useful account might show that an unreliable pipe repeatedly forces emergency deliveries, or that one area faces avoidable losses. Whether an EBU account actually improves those decisions must be tested. Even a physically accurate calculation cannot itself confer a right, repair exclusion or distribute political power.

\practice{A full reservoir serves an empty clinic tank. List three possible explanations before calling the problem physical scarcity.}{The connecting pump might be broken; the route might not have enough delivery capacity during the relevant interval; or access might be restricted by an institutional or affordability rule. The reservoir's quality and other demands also need checking. The facts given do not choose among these explanations.}

\carry{Scarcity has physical and institutional dimensions, and burdens fall unevenly. These differences matter when we turn to environmental changes that affect the conditions within which people act.}

\chapter{The condition of the planet and conditional futures}\label{ch:planet}
\question{What do environmental observations establish, and what do future pathways mean?}

An environmental statement needs a time, a place and a measurement. ``Warmer'' requires a reference period. ``Threatened'' is not the same as extinct. A projection describes what follows under specified assumptions; it does not observe the future. Keeping these distinctions clear makes the motivation stronger, because the reader can see what the evidence actually supports.

The examples here are selected assessment findings, not a complete planetary dashboard. They retain their publication dates and data periods. This edition does not present a historical assessment as a live September 2026 estimate.

\section{An observed change}
The IPCC's 2023 synthesis reports global surface temperature in 2011--2020 at 1.09 degrees Celsius above 1850--1900, with a very likely range of 0.95--1.20 degrees Celsius. The bracket describes the assessment's 5--95\% range. The same assessment attributes global warming unequivocally to human activities, principally greenhouse-gas emissions.\cite{ipcc}

This is a change in a global average over a decade. It is not the temperature in the clinic, a maximum temperature during a heatwave or a count of people exposed. An annual observation and a multi-year warming measure answer different questions. Changing the reference period without saying so would change the meaning of the number.

A global average is useful for describing a planetary change. A local decision requires further information. The clinic's water supply, building design and ability to operate during heat cannot be read from that average alone. This is an example of why a broad indicator and a local physical account can both be needed.

\section{Conditional pathways}
The IPCC also assesses contrasting greenhouse-gas scenarios. Its estimates for average warming in 2081--2100, relative to 1850--1900, include the following.\cite{ipcc}

\begin{center}\small
\begin{tabular}{@{}p{157pt}rr@{}}
\toprule
Scenario & Best estimate & Very likely range\\
\midrule
Very low emissions (SSP1-1.9) & $1.4\,^{\circ}\mathrm{C}$ & $1.0$--$1.8\,^{\circ}\mathrm{C}$\\
Intermediate (SSP2-4.5) & $2.7\,^{\circ}\mathrm{C}$ & $2.1$--$3.5\,^{\circ}\mathrm{C}$\\
Very high emissions (SSP5-8.5) & $4.4\,^{\circ}\mathrm{C}$ & $3.3$--$5.7\,^{\circ}\mathrm{C}$\\
\bottomrule
\end{tabular}
\end{center}

These are conditional climate responses to named pathways, not probabilities that society will choose each pathway. They are not a current-policy forecast. The ranges within a row describe uncertainty in the assessed response; the differences between rows also reflect different emissions assumptions. The IPCC assesses escalating risks with additional warming, with very high confidence.\cite{ipcc}

A reader should therefore resist two shortcuts. One is treating an adverse pathway as an inevitable destiny. The other is treating the existence of a lower-emissions pathway as proof that it will happen. Choices and constraints still matter. A table can clarify their consequences without making them for us.

\section{Biodiversity, materials and pollution}
In its 2019 global assessment, IPBES estimated that around one million animal and plant species were threatened with extinction. This is an assessment of threat, not a count of species already extinct. The estimate draws on assessed groups and estimated total diversity; insect threat is an important uncertainty. IPBES labels the supporting estimate established but incomplete.\cite{ipbes}

A threatened-species estimate has different units and meaning from a temperature change. It cannot be added to degrees Celsius to obtain a scientific ``damage total''. Habitat, species and ecosystem function also need more than a single worldwide count if we want to understand a particular place.

The UNEP International Resource Panel's 2024 outlook reports growth in material extraction and describes a pathway in which global extraction could rise about 60\% from 2020 to 2060 without a change of course. This is a conditional scenario, not a measured 2060 quantity. Its official overview supplies no uncertainty interval for that headline; none is invented here.\cite{resources}

Material quantity is itself incomplete as an impact measure. A tonne tells us mass. To judge consequences, we would also need to know what material, obtained where, by which process, for what use and with what residuals. The account must not acquire ecological meaning merely because it uses physical units.

Pollution adds another dimension. WHO's October 2024 assessment page describes a substantial worldwide health burden from outdoor air pollution using estimates for 2019.\cite{air} This book uses that qualitative finding without treating unlike health, climate and biodiversity measures as independent quantities that can simply be summed.

\section{Why this evidence matters to the question}
These findings make physical conditions difficult to dismiss as something outside an economy. Production, infrastructure and service depend on such conditions and can change them. Their effects are distributed unevenly, as the household example in Chapter~\ref{ch:access} illustrates conceptually.

Yet none of the assessments evaluates the EBU proposal developed here. They do not select its reference values, establish its scalar weights or test its actor policy. We should carry forward the motivation without borrowing authority for an untested remedy.

\practice{Is the difference between two scenario rows an observed temperature reduction caused by EBU?}{No. Both rows describe conditional future climate assessments. Neither includes an EBU intervention. Their contrast explains why assumptions matter, not what this project has achieved.}

\carry{Diagnosis gives us reasons to ask for better physical visibility. The next step is to specify what a proposed account would need to contribute and how it could be found inadequate.}

\chapter{From diagnosis to a research question}\label{ch:research}
\question{What would an additional physical-action account need to provide?}

Suppose the clinic's community is considering a new pump. The financial question asks whether it can be purchased and maintained within a budget. The engineering question asks whether it can deliver the required service. The access question asks who benefits and who may still be excluded. The environmental question asks what the whole represented process draws from and leaves in its surroundings.

A useful new account must improve our ability to answer some of these questions without pretending to answer all of them. It must also be worth the resources required to measure, verify and govern it. More records are not automatically better records.

\section{There is already work to build on}
Environmental accounting is not an empty field awaiting EBU. The UN's environmental-economic accounting system includes a framework for ecosystem extent, condition and services. Its physical accounts connect information about ecosystems with human activity. The framework was adopted in March 2021; that statement does not mean every possible monetary valuation has identical statistical-standard status.\cite{seea}

Environmental regulation, public provision and engineering control also already connect observations to action. Monetary policy has its own objectives and tools, as Chapter~\ref{ch:money} showed. Their purposes, spatial scales and information requirements differ. Their existence is a reason to make the proposed contribution precise.

The EBU question considered here concerns a declared action, a permitted local description, a finite before-and-after comparison and an auditable signed result. A possible contribution would lie in how those elements work together in a particular domain. This is a research proposal, not a claim of priority over every earlier accounting or allocation method.

\clearpage
\section{Seven questions for a proposed account}
Before calculating a pump-related value, we would ask:

\begin{enumerate}
\item Which physical quantities are represented, and where is the boundary?
\item Which states and actions are physically feasible?
\item What reference and scale give a change its declared meaning?
\item What information is allowed at the moment of valuation?
\item Does the calculation include the complete finite change, including shared effects?
\item Which burdens are already in the state account, and which justified residual process burdens remain?
\item What evidence connects the proposed action to what actually occurred?
\end{enumerate}

These are design requirements, not promises of success. A model can answer them clearly and still be a poor representation of the domain. Clear failure is valuable: it tells us where a claim stops instead of letting an attractive equation conceal a missing physical quantity.

For example, an account of water volume cannot by itself certify water quality. If quality is decisive, the modeller must represent it appropriately or report the account's limitation. Declaring the volume calculation exact does not repair that omission. The same reasoning applies to infrastructure condition, delayed damage and access.

\section{What would count against the proposal?}
The proposed architecture would face a mathematical failure if a supposedly exact local value omitted an affected term. It would face a measurement problem if accepted quantities could not be connected reliably to observations. It would face a practical problem if verification cost exceeded the benefit of the information. It would face an institutional problem if apparently neutral rules systematically denied essential access without a defensible justification.

These failures need different investigations. An algebraic proof can settle the first kind of claim under explicit assumptions. It cannot settle all the others. A software check may show that a formula was implemented faithfully while leaving the formula's real-world relevance unresolved.

The prospective Gaussian research programme asks a more restricted behavioural question: what happens when actors choose among feasible actions under a separately declared capacity rule? A useful study must permit adverse answers such as trapping, repeated refusal or poor recovery. The motivation for the project cannot be used to classify every outcome as success.

\section{Commitments and choices}
Transparency, contestability and the protection of essential access are commitments that guide the inquiry. They do not emerge from differentiating a potential. Who declares a reference, who can challenge a measurement and who carries the cost of a necessary action remain institutional questions.

Consider repair work that uses materials now but may improve service later. The current physical effects should be recorded under the current account. Evidence of changed capacity can justify a prospective update to the model. Responsibility for financing the repair and compensation for the worker need their own arrangements. A forecast of future usefulness should not be inserted as an immediate observed improvement.

This way of framing the project is demanding. It asks EBU to be clear enough to audit and modest enough to reject when it fails. It also leaves room for a useful result short of a complete economic system: a particular accounting method might help in one domain and fail in another.

\practice{A pump repair leaves the tank level unchanged. Has the repair achieved nothing?}{The stock measurement alone cannot answer. A repair may change equipment condition or future feasible delivery. Those require their own evidence and representation. The immediate material use can still be recorded. Missing a capacity coordinate is a model limitation, not proof that the work was worthless.}

\carry{To ask an account to preserve function, we must first explain what preserving function means. The next chapters introduce homeostasis and balance without assuming that a valuation formula already supplies a controller.}

\chapter{Homeostasis: maintaining function through change}\label{ch:homeostasis}
\question{How can a system remain functional while its components keep changing?}

The clinic's tank need not remain motionless to provide reliable service. Water enters, water leaves, demand varies and the pump starts and stops. What matters is whether the system maintains the conditions required for its function, within declared limits and disturbances.

This is a useful route into homeostasis. In physiology, homeostatic regulation concerns maintaining selected internal variables within ranges compatible with function. McFarland and colleagues' teaching framework emphasises regulated variables, continuing regulatory mechanisms and dependence on available environmental resources.\cite{homeostasis} The idea is active maintenance, not perfect constancy.

\section{A controller has to do something}
Consider a deliberately simple engineering arrangement. A level sensor measures water in a tank. A controller compares the measurement with a declared operating range. When the level falls sufficiently, it requests pumping. The motor and valve supply the physical response, provided power and source water are available. Demand and leakage disturb the level.

Each part has a role. The sensor provides information. The operating range states a criterion. The control rule turns information into a request. The pump changes the physical state. A reserve provides time to respond. A record documents what happened. Removing one role may change the behaviour of the whole arrangement.

Notice where choice entered: the control rule was explicitly specified. A scale drawn beside the tank did not start the pump. This distinction will matter later. An EBU potential can describe deviation and an EBU calculation can evaluate an action. Neither description secretly contains a pump command.

\illustration{regulation}{Schematic engineering example: a declared level-control rule connects measurement to a pump response. The tank also exchanges water with its surroundings. This is an illustration of feedback, not an EBU controller or a simulated recovery trace.}{fig:regulation}

\section{The conditions of maintaining a function}
Suppose the controller is working but the upstream supply is empty. The requested response may be impossible. Suppose the supply exists but the sensor is late. The pump may act on a state that has already changed. Suppose the pump is powerful but the delivery pipe is too narrow. The limiting component may lie elsewhere in the chain.

These are hypothetical engineering possibilities, not results of a clinical or EBU study. They explain why ``has feedback'' and ``will recover'' are different statements. A recovery claim needs a disturbance class, a specified mechanism, available resources and a definition of successful return.

A reserve also has a precise role. Stored water can sustain service during a delay, but its presence is not an unlimited guarantee. The reserve must be compared with the rate of use and duration of interruption. A tank can remain above zero and still fail a service requirement if the outlet cannot supply water fast enough.

\section{Functions can conflict}
Maintaining one measured variable may consume another resource. A pump can hold tank level while drawing down its source. A refrigerator can maintain an internal temperature while requiring electricity and rejecting heat. These examples remind us to declare a boundary wide enough for the claim being made.

Even the word ``normal'' needs care. A familiar value is not necessarily a desirable value. A preferred operating range is not automatically a measured population average. In biology, actual regulatory arrangements are more complex than our tank illustration. The analogy helps identify questions; it does not establish that economies obey the same equations as organisms.

For the clinic, function includes more than water level. Quality, timing, equipment condition, staff and access may matter. A simplified model is allowed to represent fewer things, provided its conclusions stay within that smaller world. We should say ``this model evaluates these water-stock changes'', not quietly claim that it certifies the health of the institution.

\section{Measurement and response remain separate}
It is possible to build a controller that uses a valuation signal. It is also possible to ignore the signal, to use it only for audit, or to choose randomly among actions permitted by another rule. Those are different models of behaviour.

The distinction protects the research question. If we inserted an automatic restoring mechanism and then observed return toward a reference, we would need to ask how much of the return came from that mechanism. The introductory construction does not assume such motion. We will first learn what the account measures.

\practice{The level sensor reports that the tank is below its target. Name the extra facts needed before concluding that it will refill.}{We need a response rule, an available water source, a working pump and route, enough power, relevant delay information and a feasible balance between delivery and demand. The measured discrepancy alone supplies none of these.}

\carry{Homeostasis means maintained function through a specified regulatory process. The next chapter separates that idea from equilibrium, steady state, stability and recovery, terms that are often used as if they meant the same thing.}

\chapter{Functional balance, equilibrium and living nonequilibrium}\label{ch:balance}
\question{Does balance mean that nothing moves?}

A tank receives five litres each minute and delivers five litres each minute. Its level remains constant. Water nevertheless moves through it. Looking only at the level, we might call the situation steady; looking at the flows, we would not call it inactive.

This small example prevents a large misunderstanding. A maintained condition can depend on continuing activity. It is possible to stop the activity and destroy the condition one wished to preserve.

\section{A reference is a statement; an equilibrium is a property}
A reference tells us where we want to compare a state. We might mark a tank at ten litres because that is the reference chosen for an illustrative calculation. The mark remains on the tank whether the pump is running or broken.

An equilibrium of a specified dynamic law is a state that the law leaves unchanged under the stated inputs. To decide whether ten litres is an equilibrium, we need to know what the inflow, outflow and response rule do at ten litres. The reference label cannot answer that question.

A steady state concerns selected quantities that remain constant over time. It can include throughput, as in the five-litres-per-minute example. Depending on which variables the model includes, the tank level can be a fixed point of its state equation while the larger physical apparatus keeps exchanging matter and energy. The apparent disagreement disappears when we state the model boundary.

Thermodynamic equilibrium has a different meaning: the macroscopic driving differences for spontaneous processes have equilibrated under the specified constraints. A maintained living regime commonly requires continuing exchanges and dissipation. We should not equate a healthy organism with a closed system at thermodynamic equilibrium. Lan and colleagues' analysis and experiments on bacterial sensory adaptation provide a specific example in which maintaining adaptation requires continuing energy consumption.\cite{adaptation}

The lesson for EBU is conceptual. Our dimensionless potential will not be thermodynamic free energy. Naming a mathematical field does not supply a thermodynamic law for an economy.

\section{Staying inside, staying near, and returning}
Several mathematical words describe different requirements.

\textit{Forward invariance} means that, for a declared dynamic law and allowed inputs, starting inside a specified set keeps the state inside that set. For a tank model, the set might be the interval of physically allowed levels. This says where the state cannot leave, not where it must go.

\textit{Boundedness} means that the state stays within some finite bounds. A tank can alternate between nearly empty and nearly full and still be bounded. If its capacity is finite, a bound may follow simply from the physical domain.

\textit{Stability near an equilibrium} means, roughly, that sufficiently small initial disturbances remain small under the specified dynamics. \textit{Attraction} asks for more: states in a stated region approach the equilibrium as time passes. Stability and attraction need their own mathematical definitions when used in a theorem.

\textit{Recovery} is an operational claim: after a specified disturbance, the system returns to a defined functional condition, perhaps within a required time. It depends on what counts as return and which disturbances are allowed. A system can remain bounded while never recovering.

Imagine a marker that repeatedly crosses the middle of a finite ruler. Crossing the reference does not mean settling there. The marker can visit it repeatedly without approaching rest. This is a schematic distinction, not a generated trajectory. We need a rule of motion before drawing conclusions about its future.

\section{Conservation answers another question}
If two tanks exchange water without loss, their combined quantity stays the same. Yet one can become empty while the other holds nearly everything. Conservation is an inventory statement. It does not state that each location has an appropriate amount.

Later, we will use a potential to distinguish such distributions. Even then, a low potential will describe closeness to a declared reference, not guarantee access, water quality or recovery. A poor reference can be reached exactly. A good reference can remain unreachable. The geometry and the dynamics are separate questions.

\section{A practical reading rule}
Whenever someone says a system is ``in balance'', ask which quantities, which boundary and which sense of balance they mean. Is a conserved stock unchanged? Are selected variables steady? Is an equilibrium stable? Is function being maintained by a controller? Has recovery actually been observed under a registered disturbance?

These are not pedantic distinctions. A proof of one cannot be used as evidence of all the others. They will let us understand the useful identities later without turning them into promises they do not contain.

\practice{A closed two-tank system always contains twenty litres. Is this a recovery theorem?}{No. It establishes a conserved total. A state such as $(20,0)$ can preserve that total indefinitely if the action rules leave it there. To establish recovery toward $(10,10)$, we need additional dynamics and a proof or appropriately scoped evidence.}

\carry{Before asking how a model behaves, we must specify what it contains. We now construct a small physical world with explicit stocks, units, boundaries and actions.}

\chapter{A physical world before an economic account}\label{ch:world}
\question{What must a model record so that an action has a physical meaning?}

We begin with two tanks, Ridge and Vale. Each holds the same homogeneous resource. In the illustration we call the unit a litre, written L, and assume equal water quality and fixed conditions under which the volumes add. A real application would need to justify that simplification. Continuously divisible means that fractions of a litre are allowed; it does not mean all quantities must be integers.

The model is closed during the action we evaluate: nothing enters or leaves the two-tank boundary. Transfers are idealised as lossless. We exclude pumping energy, wear and water-quality differences from this teaching world. Those omissions make it manageable, not physically complete.

\section{Stock, list and total}
Let Ridge hold fourteen litres and Vale six. A compact state description is
\begin{equation}
 x=(14,6)^{\mathsf T}\ \mathrm{L},\qquad M=14+6=20\ \mathrm{L}.
\end{equation}
The letter $x$ is a list of stocks in a fixed order. The superscript $\mathsf T$, called transpose, says to treat the displayed horizontal list as a column. It is not a power. The letter $M$ names the total represented resource.

A subscript identifies an entry: $x_1$ is Ridge and $x_2$ is Vale. A general world has finitely many cells, numbered $1$ to $n$. A cell is an accountable portion of the model, not necessarily a biological cell or a point on a geographic grid. The notation $\sum_i x_i$ means add the stocks of all its cells.

\section{A finite transfer}
Move three litres from Ridge to Vale. Write this finite quantity as $q=3\ \mathrm{L}$. The prime on $x'$ means the state after this action. Then
\begin{equation}\label{eq:transfer}
 x_i'=x_i-q,\qquad x_j'=x_j+q.
\end{equation}
Here $i$ names the source and $j$ the destination. The calculation gives $x'=(11,9)^{\mathsf T}\ \mathrm{L}$. The total remains twenty litres because subtracting three at one place and adding three at the other cancels in the sum:
\begin{equation}
 \sum_i x_i'=\sum_i x_i=M.
\end{equation}
This is a conditional conservation identity. It assumes the stated lossless internal transfer. It is useful because a broken total points to a missing crossing, an incorrect action or inconsistent measurements. It does not tell us whether $(11,9)$ is desirable.

A feasible action cannot withdraw more than the source contains or overflow a declared physical storage capacity. If both tanks have a physical capacity of twenty litres, the example satisfies those restrictions. Route limits, permissions to use the pipe and other constraints must be declared if they are part of the model. They are not supplied by subtraction.

\illustration{transfer}{Schematic with exact arithmetic: a three-litre lossless transfer changes $(14,6)$ to $(11,9)$ and preserves the twenty-litre total. Values are invented for teaching; no model was executed.}{fig:transfer}

\section{Rate is not quantity}
A rate says how much moves per unit time. If a constant rate is $r=2\ \mathrm{L/min}$ over an interval $\Delta t=1.5\ \mathrm{min}$, then
\begin{equation}\label{eq:rate}
 q=r\Delta t=3\ \mathrm{L}.
\end{equation}
The minutes cancel when the units multiply. Once $q$ has been formed, we do not multiply it by time again. If the rate varies, the total quantity must come from the measured or declared accumulation over the interval; a single endpoint rate generally does not suffice.

Some historical EBU sources use the symbol $q$ for a rate. In this edition $q$ always denotes a finite quantity. The convention is explicit because an unnoticed change of meaning can produce a correct-looking equation with incorrect units.

\section{What changes when a boundary is open?}
Suppose instead that three litres leave Ridge and only 2.7 arrive at Vale. The represented stocks fall by 0.3 litres. That amount needs a named physical destination or boundary account, such as leakage outside the measured tanks. Calling it ``loss'' does not explain where the matter went.

This physical bookkeeping does not yet determine an EBU process charge. We first record the transfer and its boundary crossings; later we ask which effects are represented in the potential and whether any justified residual burden remains. A loss ledger and a valued potential factor are distinct records.

Likewise, repairing the pipe may change its physical capacity while leaving current water volume unchanged. A stock-only list cannot represent that improvement. We would need evidence of the changed property and a justified prospective model update. Adding a hopeful future story to the present stock difference would not fix the missing coordinate.

\practice{Starting from $(14,6)$ litres, calculate the endpoint of a six-litre transfer. Is conserving twenty litres enough to show improvement?}{The endpoint is $(8,12)$ litres and the total is twenty. Conservation alone cannot compare it with the start. We need a declared reference and measure, introduced in Chapters~\ref{ch:reference}--\ref{ch:potential}.}

\carry{An action is a physical change before it is a value. The next chapter places that change among feasibility, information, permission and actor choice.}

\chapter{Actions, actors and local information}\label{ch:actors}
\question{Who acts, what is evaluated, and what information is allowed?}

The transfer has a source, a destination and a quantity. An actor can propose it. These are separate facts. A person may make several proposals; several people may cooperate on one process; an automatic device may perform the physical operation. The mathematical description of the action does not by itself assign ownership or responsibility.

For our simple world, write an action as $a=(i,j,q)$. This names a lossless transfer of the finite quantity $q$ from cell $i$ to cell $j$. The notation is a compact instruction for a represented transformation. Its physical validity depends on the world and constraints already declared.

\section{Seven layers of one event}
The working sequence is:

\begin{center}\small
physical state and proposed action\\$\downarrow$\\
physical feasibility\\$\downarrow$\\
valuation\\$\downarrow$\\
separately declared permission or account policy\\$\downarrow$\\
actor choice\\$\downarrow$\\
execution\\$\downarrow$\\
audit
\end{center}

Each layer answers a distinct question. Feasibility asks whether the represented transformation can occur within physical constraints. Valuation asks what its declared potential difference would be. A separate policy may determine which feasible actions are authorised or affordable. The actor's rule then selects an action or declines to act. Execution changes the physical state. Audit checks what occurred against the declared record.

In a real institution, some legal or safety permissions may be checked earlier as well. The conceptual separation remains: a favourable value is not physical capability or legal permission. An unfavourable value does not decide a person's entitlement to essential service.

\section{What local means}
To evaluate a simple two-cell transfer, the account will need the two stocks, their reference parameters, the action quantity and the applicable process-burden declaration. It need not read an unrelated tank whose potential term is unchanged. Later we prove this cancellation; at present it is a requirement to justify, not a licence to ignore affected effects.

Locality follows dependencies, not mere geographic distance. A long pipe can connect two model cells. A nearby third cell might be irrelevant to one action, while a farther cell could be relevant if the declared potential couples it to the action. The information boundary has to match the physical and mathematical description.

The valuation should not quietly consult a simulated future, a global ranking, a person's wealth or an eventual study outcome. A planning system may use different information for its own task, but that is a separate mechanism. Mixing its forecast into an allegedly present local value changes the claim.

\section{Several actions can occur together}
Suppose two actors both propose to deliver water into Vale from separate tanks. The destination changes once under the combined delivery. The pair therefore needs a coherent joint endpoint and a joint feasibility check. Each action being feasible alone does not ensure that their combination respects a shared capacity.

There is no rule in this introductory world that only one action may be accepted per micro-step. Simultaneous groups remain in scope. We will learn in Chapter~\ref{ch:groups} why simply adding independent values evaluated from the same start can miscount the joint change.

If one delivery truly occurs after another, the second begins from the state the first left. That is a sequential description. It is not equivalent to pretending that both began from the original state. Nor may we invent a physical order merely because a computer happens to list one actor first.

\section{Keep transformation chains visible}
The clinic's real pump belongs to a chain: an energy source supplies electricity, equipment converts it into mechanical work, and a pipe delivers water. When a model represents those processes, their separate physical transformations and boundaries should remain traceable. One purchase need not become one giant physical action that hides all the intermediate effects.

Conversely, several links do not justify charging an extra burden merely because there are several links. An effect must have a physical meaning and be counted once. A regional audit must not count an internal transfer again as new external supply.

\section{An actor rule is an additional choice}
A researcher could specify random choice, a service priority, a planner or a value-maximising policy. Those policies can produce different outcomes even with the same value equation. EBU calculation does not choose among them.

The prospective Gaussian programme considers EBU-blind candidate formation and random choice after separate feasibility and affordability checks. It does not secretly rank by EBU or add a restoring force. This is a proposed behavioural study, not a prescription for patient care or an implemented economy.

\practice{An action has a positive calculated value but would overflow the destination. Should the actor treat the value as permission?}{No. The action fails the physical constraints of the declared world. Valuation cannot make an invalid endpoint feasible. Equally, a feasible negative-valued action raises a separate policy question; its sign alone is not a universal prohibition.}

\carry{We now know what an action and an actor are. To evaluate the action, we need to declare the ruler by which a represented change counts as nearer to or farther from a reference.}

\chapter{Why a reference is necessary}\label{ch:reference}
\question{Relative to what does a measured change count as improvement?}

Three litres moved. That statement describes a physical quantity. It does not yet tell us whether the move improved the represented condition. A full tank may have needed relief, or it may have held a necessary reserve. The destination may have needed water, or it may already have been too full. A direction alone does not supply the criterion.

A reference makes the comparison explicit. It says which stock at each cell anchors the ruler. We write the reference at cell $i$ as $x_i^\star$, pronounced ``x i star''. The star is a label for the declared reference, not multiplication and not a proof that the value is optimal.

\section{A compatible reference}
For the two-tank illustration, choose ten litres at Ridge and ten at Vale. The total reference is twenty litres, matching the total available water. These are teaching choices, not calibrated operating recommendations for a real clinic.

For a closed world with fixed total $M$, a physically compatible reference must satisfy
\begin{equation}\label{eq:compatibility}
 x_i^\star\geq0,\qquad \sum_i x_i^\star=M.
\end{equation}
If a cell has a declared physical storage capacity $K_i$, it must also satisfy $x_i^\star\leq K_i$. All these quantities use the resource's stock unit. The conditions ensure that the reference does not demand negative stock, an unavailable total or a physically impossible volume at one cell.

Compatibility is necessary for reaching the reference in this domain, but it is not enough. A disconnected pipe network might prevent water reaching a tank. A directed route might allow transfer only the wrong way. A restricted action menu might omit the required quantities. A minimum on paper cannot create a missing route.

A reference need not assign equal stocks to all cells. Different roles may justify different values. This book uses equal values so that the arithmetic is easy to follow. Equality here is an illustrative symmetry, not a theorem about fair allocation.

\section{A scale says how to read a displacement}
We also declare a positive scale $\sigma_i$, pronounced ``sigma i'', in the same stock unit. The scale tells us how large a displacement is in the ruler's own terms.

A two-litre displacement is one scale unit when $\sigma_i=2\ \mathrm{L}$, but two scale units when $\sigma_i=1\ \mathrm{L}$. The water did not change when the scale changed. The representation changed. Any scientific or institutional use of that difference needs a justification for the scale.

In our recurring example, both scales are two litres. Ridge's displacement from reference is $14-10=4\ \mathrm{L}$; Vale's is $6-10=-4\ \mathrm{L}$. Dividing each by its scale gives $+2$ and $-2$, dimensionless numbers. The sign says which side of the reference the stock occupies.

The scale is not automatically a tolerance, a safety margin, a measurement error or an observed standard deviation. It may be given one of those interpretations in a separately justified model, but its notation does not establish that interpretation. A positive scale is required because dividing by zero would leave the ruler undefined.

\section{Who chooses the ruler?}
For an application, the modeller needs a basis for the reference and scale: domain evidence, the intended function, a measurement method, uncertainty and a procedure for review. If the declaration affects rights or access, governance must also be explicit.

A measured historical mean is not automatically a desirable reference. A system might have spent years in a deficient condition. Nor does a preferred reference automatically describe what its dynamics will maintain. Chapter~\ref{ch:balance} already separated a declared mark from an equilibrium.

Versions matter. An action evaluated under one ruler must not be quietly re-evaluated under another because the second yields a preferred sign. If evidence justifies changing the model, old records retain their original declaration and the new version applies prospectively under declared rules. A reference is part of the account's identity.

\section{What this choice leaves open}
The simple Gaussian construction will penalise displacement in both directions. That makes a transparent teaching ruler, but real systems may have asymmetric risks, viable ranges or thresholds. Such structures need other justified potentials. Introducing this quadratic does not invalidate the historical threshold family or declare nature universally quadratic.

Normalising several quantities also does not establish an ethical exchange rate between them. Dividing water by a water scale and medicine by a medicine scale does not prove that their scores may be combined into one universal welfare number.

\practice{There are twenty litres, but the proposed reference is twelve litres in each of two tanks. What is wrong?}{The reference requires twenty-four litres. Its zero-deviation state is unavailable in the closed twenty-litre domain. A calculation can still evaluate displacements from it, but describing zero deviation as a physically compatible target would be false.}

\carry{With a declared reference and positive scale, we can build a simple ruler. The next chapter explains every operation before giving that ruler its Gaussian name.}

\chapter{The local Gaussian ruler}\label{ch:gaussian}
\question{How can distance from a reference be measured transparently?}

Start with the displacement: stock minus reference. Divide by the positive scale to express that displacement in comparable local scale units. Square the result so that equal-sized displacements above and below the reference contribute equally. Finally, multiply by one half, a convenient convention that simplifies the derivative later.

These four operations define a local potential:
\begin{equation}\label{eq:local}
 V_i(x_i)=\frac12\left(\frac{x_i-x_i^\star}{\sigma_i}\right)^2,
 \qquad \sigma_i>0.
\end{equation}
The subscript $i$ identifies a cell. The function $V_i$ turns its stock into a declared deviation score. Both numerator and denominator have stock units, so their ratio and its square are dimensionless. In this illustrative account, potential and EBU values therefore use a dimensionless account unit. They are not litres, joules or currency.

\section{Calculating before interpreting}
For a ten-litre reference and two-litre scale, a stock of fourteen litres gives
\begin{equation}
 V_i(14)=\frac12\left(\frac{14-10}{2}\right)^2
 =\frac12(2)^2=2.
\end{equation}
A stock of six litres gives the same value: the standardised displacement is $-2$, whose square is four. At the reference, the displacement and potential are both zero.

For a stock of twelve litres, the value is $\tfrac12(1)^2=0.5$. Doubling a displacement from two litres to four multiplies its squared contribution by four. This is a property of the chosen ruler. It is not an empirical assertion that every physical harm grows quadratically.

If the scale were four litres instead of two, a four-litre displacement would contribute $0.5$ instead of $2$. A wider scale makes the same displacement count less. This is why scale selection cannot be hidden in implementation or adjusted after seeing an outcome.

\illustration{quadratic}{Mathematically derived: $V_i=\tfrac12((x_i-10)/\sigma_i)^2$ for scales of two and four litres. Dashed and solid curves distinguish the scales without relying on colour. The curves are reference geometry, not measured fluctuations or simulated motion.}{fig:quadratic}

\section{Why the word Gaussian?}
A normal, or Gaussian, density has a familiar bell shape. In its standard mathematical definition, a location parameter sets its centre and a positive scale sets its width.\cite{normal} Using our symbols, a reference density on the real line would be
\begin{equation}\label{eq:normal}
 p_i(u)=\frac{1}{\sigma_i\sqrt{2\pi}}
 \exp\!\left[-\frac12\left(\frac{u-x_i^\star}{\sigma_i}\right)^2\right].
\end{equation}
Here $u$ is a possible stock-coordinate value in the mathematical reference, $\pi$ is the circle constant and $\exp$ is the exponential function. A density has inverse-stock units; probabilities would come from integrating it over intervals under an adopted probability model.

Compare the density at $u$ with its value at the centre. The identical normalising factors cancel. The negative natural logarithm of that ratio is
\begin{equation}\label{eq:log}
 -\log\frac{p_i(u)}{p_i(x_i^\star)}
 =\frac12\left(\frac{u-x_i^\star}{\sigma_i}\right)^2.
\end{equation}
The logarithm reverses the exponential, and the minus sign makes the score grow away from the centre. The density ratio is dimensionless, so taking its logarithm is consistent with the units. This relative-log-density identity explains the name ``Gaussian ruler''.

A reader need not use probabilities to use the quadratic definition. The bell shape explains an established mathematical connection. It does not provide evidence that the actual stocks or future trajectories follow a normal distribution.

\section{The physical domain does not become Gaussian}
The normal reference on the real line assigns density even to negative values. Our physical tanks do not permit negative stocks. They also have a fixed combined total, which couples their possible states: if one stock rises in a lossless transfer, another falls.

Consequently, a sum of local quadratic terms does not establish independent normally distributed physical stocks. We use the reference geometry on a separately declared physical domain. A probability law for actual states would require additional assumptions about the process and would need separate justification.

The name also has nothing to do with Gauss's divergence theorem, a theorem relating a field's flux through a boundary to its divergence inside. Neither that theorem nor a bell curve supplies a law of economic recovery. Here the connection is exactly the relative logarithm in equation~\eqref{eq:log}.

A final caution concerns similar-looking functions. The score $1-\exp(-V_i)$ is bounded, whereas $V_i$ grows without a bound on the full real line. Their slopes and finite differences differ. Replacing one with the other changes the model; it is not cosmetic notation.

\practice{With reference ten litres and scale two litres, what are the scores at eight, ten and twelve litres?}{They are $0.5$, $0$ and $0.5$. The equality on either side expresses symmetry of the ruler. It does not say that the corresponding physical conditions have equal consequences in every real application.}

\carry{We have defined a local score and stated its limits. Next we add such scores to describe a joint state, while preserving the possibility of exact local evaluation.}

\chapter{From local references to a potential}\label{ch:potential}
\question{How can local measurements describe a joint state without every actor inspecting the whole world?}

Ridge contributes a local score of two and Vale contributes two. To describe their joint state, declare the total potential to be the sum of the local terms:
\begin{equation}\label{eq:potential}
 V(x)=\sum_{i=1}^{n}V_i(x_i).
\end{equation}
The summation sign means ``add for every cell from one to $n$''. The list $x$ contains all stocks. In the two-tank illustration, $V(14,6)=2+2=4$.

The addition is part of the model. It says these local contributions enter one compatible account. It does not prove that all ecological and human consequences can be represented by one scalar. In this chapter, all cells contain one homogeneous resource and use the explicitly declared ruler.

\section{Same total, different potential}
Compare three states, each containing twenty litres. At $(14,6)$, the potential is four. At $(10,10)$, both displacements vanish and the potential is zero. At $(18,2)$, the standardised displacements are $+4$ and $-4$, giving $8+8=16$.

The physical total cannot distinguish these distributions. The potential can, because it measures location-specific displacement from the reference. Neither operation creates or destroys water. The potential is not a conserved stock.

This also shows why the word ``burden'' has to be qualified. In this model it means represented squared deviation. It does not mean an independent observation that one state caused sixteen units of real-world harm. Interpretation requires a justified link between the chosen reference geometry and the physical function of interest.

\section{Why unchanged terms cancel}
Imagine extending the illustration with two other tanks. They hold twelve and eight litres, each with reference ten and scale two. The four-cell world now contains forty litres, exactly matching the four references. Each of the extra tanks contributes $0.5$.

Our original three-litre Ridge-to-Vale transfer leaves those two extra tanks unchanged. The full potential is five before the transfer and $1.25$ afterwards:
\begin{align*}
 V_{\rm before}&=2+2+0.5+0.5=5,\\
 V_{\rm after}&=0.125+0.125+0.5+0.5=1.25.
\end{align*}
The reduction is $3.75$. Looking only at Ridge and Vale gives $4-0.25=3.75$ as well. The other terms were not approximated; they cancelled exactly.

Let $A$ be the set of all cells whose terms can change in the action, and let $V_A$ be their sum. If $x'$ agrees with $x$ outside $A$, then
\begin{equation}\label{eq:locality}
 V(x)-V(x')=V_A(x)-V_A(x').
\end{equation}
To prove it, split each full sum into affected and unaffected terms. Every unaffected term has the same value on both sides and subtracts to zero. This is the separable local-to-global identity in the historical foundation, expressed here for the Gaussian model.\cite{foundation}

It is useful because an exact local difference need not require a whole-world scan. The proof still depends on one fixed potential, correct action support and unchanged terms outside it. It cannot certify that the model includes all relevant real effects.

\section{If the potential has coupled factors}
A later model might include a factor depending on two stocks together. Suppose, as a mathematical extension, it adds
\begin{equation*}
 W(x_2,x_3)=\gamma(x_2-x_3)^2,
\end{equation*}
where $\gamma$ has inverse-square-stock units so that $W$ remains dimensionless. Changing $x_2$ can change $W$ even when $x_3$ stays fixed. Reading only the cell-two term would then omit a changed contribution.

The correct affected support includes every potential factor touched by the action and the variables needed to evaluate those factors. Cancellation still works for genuinely unchanged factors. ``Local'' is therefore a dependency claim that must be checked against the actual function.

Physical topology, potential dependencies and action interaction are distinct. A pipe describes an allowed physical transfer. A potential factor describes which variables the ruler couples. Two actions can interact in valuation by changing the same coordinate even when the potential itself is a sum of independent cell terms. Chapter~\ref{ch:groups} makes that last point concrete.

\section{The global audit is not the actor's policy}
Writing $V$ is convenient for proving identities and reviewing the complete account. It does not require actors to rank actions by a global score. A theorem can mention a full sum while a local calculation uses only affected terms.

Nor does local evaluation prove decentralised governance. Someone still declares references, certifies inputs and resolves disputes. Distributed calculation and institutional authority are different matters. Mathematics can reduce the required information without choosing a constitution.

\practice{An unchanged tank contributes seven units of potential before and after an action. How much does it contribute to that action's value?}{Zero: $7-7=0$. Its current condition may matter for other questions. It contributes nothing to this finite difference only because the action leaves the term unchanged under the same potential version.}

\carry{A potential describes states. To understand the initial direction of a change, we next examine its local rate of change: the marginal.}

\chapter{Marginal field: what a small change means}\label{ch:marginal}
\question{What does the local slope say about a small transfer?}

Ridge is above reference and Vale below it. Removing a little from Ridge and adding a little to Vale reduces both contributions. We want a number for this immediate sensitivity, while remembering that a real delivery is finite.

A derivative is the limiting ratio of a change in potential to an increasingly small change in stock. It describes the starting slope, not the entire finite difference.

\section{Finding the slope without guessing}
Write $d=x_i-x_i^\star$ for the displacement in one cell and let $h$ be a small stock increment. Expanding the square gives
\begin{equation*}
 V_i(x_i+h)-V_i(x_i)
 =\frac{d}{\sigma_i^2}h+\frac{h^2}{2\sigma_i^2}.
\end{equation*}
Divide by $h\ne0$. The ratio is $d/\sigma_i^2+h/(2\sigma_i^2)$. As $h$ tends to zero, the second term tends to zero. The derivative, called the local marginal burden here, is therefore
\begin{equation}\label{eq:mu}
 \mu_i=\frac{x_i-x_i^\star}{\sigma_i^2}.
\end{equation}
The Greek letter $\mu$, pronounced ``mu'', names the slope. A dimensionless potential differentiated with respect to litres has units $\mathrm{L}^{-1}$. A marginal is neither a stock nor a completed action value.

For Ridge, $\mu_1=(14-10)/4=1\ \mathrm{L}^{-1}$. For Vale, $\mu_2=(6-10)/4=-1\ \mathrm{L}^{-1}$. A negative slope means that adding a small amount lowers the local potential. A positive slope means that adding a small amount raises it. At the reference the slope is zero.

For an arithmetic check, add $0.1$ litre to Vale while examining its local term alone. The linear prediction is $-1(0.1)=-0.1$. The exact change is $-0.1+0.1^2/8=-0.09875$. The local prediction is close for this small increment but is not identical. Examining one term this way is a sensitivity calculation, not a proposal to create water in the closed world.

\section{A direction combines two slopes}
A conservative transfer removes stock at one endpoint while adding it at another. Let $e_i$ mean a list with one in position $i$ and zero elsewhere. For edge $e$ from $i$ to $j$, define
\begin{equation}\label{eq:edge}
 S_e=-e_i+e_j,\qquad \delta x=qS_e.
\end{equation}
The direction list $S_e$ is dimensionless. Multiplying it by the stock quantity $q$ produces the increment $\delta x$: negative at the source, positive at the destination. For two cells it is simply $q(-1,+1)^{\mathsf T}$.

The list of local marginals is written $\mu$. Its dot product with a change list means multiply matching entries and add. Thus
\begin{equation*}
 \mu^{\mathsf T}\delta x=\mu_i(-q)+\mu_jq.
\end{equation*}
The transpose arranges this multiplication and addition. The result is the first-order potential change. For positive value to mean reduction, define the directional field
\begin{equation}\label{eq:field}
 f_e=-\mu^{\mathsf T}S_e=\mu_i-\mu_j.
\end{equation}
It has inverse-stock units. Multiplying by a small quantity gives a first-order field value in the potential's units.

In the recurring example, $f_e=1-(-1)=2\ \mathrm{L}^{-1}$ from Ridge to Vale. Reversing direction gives $-2\ \mathrm{L}^{-1}$. When both tanks are at reference, both marginals and the initial field are zero. These signs describe local geometry, not moral categories.

\section{A field is not a command}
Here ``field'' means local marginal information for states and directions, not an electromagnetic field. The project literature sometimes calls $f_e$ a force, but it is not measured in newtons and does not push an actor.

A flow requires an additional rate law, including feasibility and time. Selection requires an actor policy. Differentiation provides neither. An actor may decline a positive-valued option.

\practice{With $\mu_i=0.5\ \mathrm{L}^{-1}$ and $\mu_j=0.5\ \mathrm{L}^{-1}$, what is the initial directional field?}{It is zero. The first-order changes at the endpoints cancel. This does not make an arbitrary finite transfer neutral: curvature can still increase the joint potential, as the next chapter shows.}

\carry{A delivery is finite. Its value requires the entire endpoint change.}

\chapter{A real action is finite}\label{ch:finite}
\question{Why can a favourable direction become unfavourable when taken too far?}

Our starting field from Ridge to Vale is positive. That does not mean sending all of Ridge's water is beneficial under the declared ruler. During a transfer, Ridge moves down toward its reference and Vale moves up toward its reference. After both reach it, further movement takes them away again.

We need the endpoint, not just the first slope. For a proposed increment $\delta x$, first check physical feasibility and form $x+\delta x$. In the ideal lossless model with no separate process burden, define the field value as the potential before minus the potential after:
\begin{equation}\label{eq:finite-def}
 E(\delta x)=V(x)-V(x+\delta x).
\end{equation}
The sign convention is deliberate. A reduction in represented deviation gives a positive result. $E$ is in the same dimensionless account units as $V$.

\section{The whole correction is calculable}
For one cell, expansion of the square is exact:
\begin{equation*}
 (x_i+\delta x_i-x_i^\star)^2
 =(x_i-x_i^\star)^2+2(x_i-x_i^\star)\delta x_i+(\delta x_i)^2.
\end{equation*}
Divide by $2\sigma_i^2$ and add across cells. The first terms reproduce $V(x)$; the middle terms give $\sum_i\mu_i\delta x_i$; the remaining terms are the curvature contribution.

A compact way to write the last sum uses $H$, the diagonal curvature matrix:
\begin{equation}\label{eq:hessian}
 H=\operatorname{diag}(1/\sigma_1^2,\ldots,1/\sigma_n^2).
\end{equation}
``Diagonal'' means the listed entries lie on the diagonal and the other entries are zero. Consequently $\delta x^{\mathsf T}H\delta x$ is simply $\sum_i(\delta x_i)^2/\sigma_i^2$. No prior matrix knowledge is needed to evaluate it. Each entry of $H$ has inverse-square-stock units.

Subtracting the expansion from the starting potential gives
\begin{equation}\label{eq:finite}
 \boxed{E(\delta x)=-\mu(x)^{\mathsf T}\delta x
 -\frac12\delta x^{\mathsf T}H\delta x.}
\end{equation}
The first term is the initial slope prediction. The second subtracts the finite curvature correction. There are no omitted higher-order terms for this fixed quadratic geometry. ``Exact'' here describes algebra, not a guarantee of exact physical measurements or floating-point arithmetic.\cite{gaussian,foundation}

Because all scales are positive, the curvature sum is positive for every nonzero increment. The linear prediction therefore overstates the finite field reduction. This is the quadratic instance of the convexity argument in the historical foundation. We have derived the claim for this model; we have not assumed a recovery law.

\section{The transfer formula}
For a transfer of finite quantity $q$, the source change is $-q$ and the destination change is $+q$. Substitution gives
\begin{equation}\label{eq:edge-finite}
 E_e(q)=qf_e-\frac{q^2}{2}
 \left(\frac1{\sigma_i^2}+\frac1{\sigma_j^2}\right).
\end{equation}
The units confirm the calculation: stock times inverse stock in the first term, and stock squared times inverse stock squared in the second. Both are dimensionless. The formula assumes a fixed reference and scales, the stated lossless action, a feasible endpoint and zero separate process burden.

In our example $f_e=2\ \mathrm{L}^{-1}$ and both scales are two litres. For a numerical quantity $q$ expressed in litres, the value is $2q-q^2/4$. The table compares alternatives from the same starting state; its rows are not consecutive model steps.

\begin{center}\small
\begin{tabular}{@{}rrrr@{}}
\toprule
$q$ (L) & Endpoint (L) & $V$ after & $E$\\
\midrule
0 & $(14,6)$ & 4 & 0\\
1 & $(13,7)$ & 2.25 & 1.75\\
3 & $(11,9)$ & 0.25 & 3.75\\
4 & $(10,10)$ & 0 & 4\\
6 & $(8,12)$ & 1 & 3\\
8 & $(6,14)$ & 4 & 0\\
9 & $(5,15)$ & 6.25 & $-2.25$\\
\bottomrule
\end{tabular}
\end{center}

The three-litre action gives $6-9/4=3.75$, rather than the linear prediction of six. The nine-litre action is feasible with the declared capacities but has a negative value. Passing the minimum is not immediately the same as having a negative total value: the six-litre action has overshot the reference yet still ends closer than it began.

\illustration{finite_value}{Mathematically derived: alternative finite quantities from the same $(14,6)$-litre baseline. The horizontal axis is quantity, not time. The dashed initial prediction exceeds the exact quadratic value. A maximum in the graph does not select an actor policy.}{fig:finite}

\section{At the reference}
At $x=x^\star$, every marginal is zero. Equation~\eqref{eq:finite} becomes $E=-\tfrac12\delta x^{\mathsf T}H\delta x$, strictly negative for any nonzero net increment. A nonempty group of opposing actions can have zero net increment; that is a different case. In the ideal cost-free account its field value is zero, not negative merely because several actions were listed.

\practice{Starting at $(10,10)$ litres, move one litre from Ridge to Vale. Calculate the exact field value.}{The endpoint is $(9,11)$. Each cell contributes $\tfrac12(1/2)^2=0.125$, so $E=0-0.25=-0.25$. The zero initial slope did not make the finite action free of represented deviation.}

\carry{The finite endpoint difference is the core object. We now put it into the general EBU equation and explain the process-burden boundary that the ideal example has deliberately set to zero.}

\chapter{The exact EBU equation}\label{ch:ebu}
\question{What precisely is calculated for a declared physical action?}

In words, the account is: the affected potential before the action, minus the affected potential after the action, minus any separately declared residual physical process burden that has not already been counted.

The general notation is
\begin{equation}\label{eq:ebu}
 \boxed{\Delta E_a=V_A(z)-V_A(T_a(z))-C_a.}
\end{equation}
This is the book's central finite account. Its compactness should not hide the declarations needed to use it.\cite{foundation,gaussian}

\section{Read every symbol}
The subscript $a$ identifies the action. The symbol $z$ names the frozen action baseline: the state against which the action is evaluated. In the undriven tank example, $z=x$. If an external event is declared to occur before the action, $z$ is the state after that event. It is not a future outcome that the actor has secretly inspected.

The expression $T_a(z)$ names the state produced by the declared action transformation. Here $T$ means transformation, not the superscript transpose. For a lossless edge, $T_a(z)=z+qS_e$. For a more complicated process, its actual map must be specified; it cannot be inferred from the value equation.

The set $A$ contains the complete affected potential support. In the separable two-tank model it contains the source and destination. In a model with coupled factors it must include every changed factor, as Chapter~\ref{ch:potential} explained. $V_A$ sums those terms. Unchanged factors cancel, so the affected difference equals the full model difference under the cancellation assumptions.

The term $C_a$ is a declared nonnegative physical action-process burden, measured in the same account unit as the potential, subject to the no-double-count boundary. It is zero for a genuine no-action event. The symbol $\Delta E_a$ is the signed action value; the delta announces a change-related account, not a stock of physical material.

For our Gaussian ruler, $V_A$, $C_a$ and $\Delta E_a$ are dimensionless account quantities. In another calibrated potential, all three must share that potential's unit. Neither the letter E nor the project's name makes the result energy in joules.

\section{A complete small account}
Return to the three-litre delivery. Its teaching record can be read without software:

\begin{center}\small
\begin{tabular}{@{}p{107pt}p{239pt}@{}}
\toprule
Record field & Declared illustration\\
\midrule
Boundary & Two homogeneous water stocks; twenty litres total\\
Ruler & Fixed references $(10,10)$ L; scales $(2,2)$ L\\
Physical limits & Nonnegative stocks; each tank capacity twenty litres\\
Baseline & $z=(14,6)^{\mathsf T}$ L; no intervening drive\\
Action & Ridge to Vale, finite quantity three litres\\
Endpoint & $T_a(z)=(11,9)^{\mathsf T}$ L\\
Potential & $4$ before; $0.25$ after\\
Process burden & $C_a=0$ in this idealised teaching world\\
Result & $\Delta E_a=4-0.25-0=+3.75$\\
Evidence status & Exact illustrative arithmetic; no measured execution\\
\bottomrule
\end{tabular}
\end{center}

The physical total is twenty litres both before and after. The potential falls by $3.75$. The two facts are compatible because water and represented deviation are different quantities. If the action were only proposed, this would be a prospective value, not evidence that three litres had actually arrived.

A real record would need measurement uncertainty, observation times, action identity, model version and occurrence evidence. It would also need an explicitly declared policy for authorisation and settlement. Naming an actor in the record would not identify an employer, assign personal liability or decide compensation.

\section{What belongs in process burden?}
The ideal example sets $C_a=0$ because it deliberately contains only lossless stock transfers with no separately modelled residual process burden. It does not claim that real pumping, transport or maintenance is free.

A later nonzero term would need an identified physical process, compatible calibration and an audit showing that the same contribution has not already been represented through the transition, affected state coordinates or finite potential difference. Source drawdown already valued in $V_A$ cannot be charged again under another name. Money prices, wages, inconvenience and fraud penalties do not silently become physical EBU terms.

An unvalued physical ledger also does not automatically justify a process charge. The historical foundation contains a partially represented dissipation problem, and the current reconciliation leaves a loss-ledger interpretation unresolved. This edition does not settle it by choosing a convenient loss coefficient. Lossy worlds still require explicit physical carriers, sinks or boundary accounts and a separately justified valuation boundary.

If a future state description represents a physical effect previously treated elsewhere, the two declarations must be reconciled so it is counted once. ``Residual'' is not a licence to insert whatever cost makes a preferred outcome appear.

\section{Signs, drive and verification}
A positive result means the declared potential reduction exceeds the declared process burden. A zero result means the net account is zero. A negative result means the represented account has worsened after including that burden. None of these signs measures a person's worth or automatically grants or denies essential service.

If external drive changes $x$ into $z$ before an actor acts, compare $z$ with $T_a(z)$. For no action, $T_0(z)=z$ and $C_0=0$, so the actor's result is zero. Using the earlier $x$ instead would mix the external event's potential change into the actor's account. The shared-baseline discipline removes that immediate misattribution. It does not solve every question about later causal effects.\cite{foundation}

For a pre-action quote, the rule, baseline conditions and permitted quantity domain must be declared before execution. Audit can verify occurrence and the relevant measured quantity under those rules; it cannot invent a more favourable ruler afterwards. A stale baseline, overexecution or failed sensor needs its declared failure treatment. Exact algebra cannot certify an unverified event.

\practice{An action reduces represented potential by five. Someone proposes an extra charge of five for exactly the same represented source drawdown. What must the reviewer do?}{Reject the duplicate term. The effect is already inside the finite difference. A genuinely different residual physical process would require its own declaration and evidence; it cannot be established by relabelling the same five.}

\carry{The equation measures one coherent physical change. When actions share a state, we need to understand the group endpoint before assigning any individual account lines.}

\chapter{One action, several actions and the path between states}\label{ch:groups}
\question{Why does a simultaneous group need one coherent before-and-after account?}

A destination can receive water from more than one source. Each delivery can be measured, yet the value of the pair can differ from the sum of the values calculated for each delivery alone. The difficulty is not necessarily uncertainty about who delivered what. It can arise from the curvature of the declared ruler.

To see it clearly, we use a new three-tank illustration. All references are ten litres and all scales are two litres. The common starting state is $(14,12,4)$ litres, so the total and the reference total are both thirty litres. Each tank has a physical capacity of twenty litres. Action $a$ sends two litres from the first tank to the third. Action $b$ sends two from the second to the third. Both the single actions and their group satisfy these declared constraints.

\section{Calculate the group first}
The initial potential is
\begin{equation*}
 V(14,12,4)=\frac{4^2+2^2+(-6)^2}{8}=7.
\end{equation*}
If $a$ happens alone, the endpoint is $(12,12,6)$ and the potential is three, giving value four. If $b$ happens alone, the endpoint is $(14,10,6)$ and the potential is four, giving value three. These are two alternative singleton calculations against the same start.

If both actions happen, their joint endpoint is $(12,10,8)$, with potential one. The exact group field value is $7-1=6$, not $4+3=7$. The singleton sum has claimed one unit too much. Both singleton calculations included the earlier, steeper portion of the destination's relief.

\illustration{group}{Schematic with exact arithmetic: two two-litre actions share one destination. Singleton alternatives from $(14,12,4)$ have values four and three; their joint endpoint has value six. No time evolution or experimental outcome is plotted.}{fig:group}

This example uses zero separate process burden. Subtracting correctly declared costs from both sides would not by itself repair a duplicated field contribution. The joint field difference has to be right first.

\section{The algebra of shared change}
For fixed additive actions, let $\delta x_a$ be the increment of child $a$ and define
\begin{equation}\label{eq:group}
 \delta x_G=\sum_{a\in G}\delta x_a,\qquad
 E_G=V(z)-V(z+\delta x_G).
\end{equation}
The set $G$ is the group. All increments use the same frozen baseline $z$, the same fixed potential and the declared joint action map. Additivity is an assumption about those physical increments, not a universal property of arbitrary processes.

Expanding the quadratic for two actions gives
\begin{equation}\label{eq:cross}
 E_{\{a,b\}}=E_a^{(0)}+E_b^{(0)}
 -\delta x_a^{\mathsf T}H\delta x_b.
\end{equation}
The superscript $(0)$ means a singleton value evaluated against the common start, not the state left by another action. The cross term measures the shared curvature contribution. It is dimensionless, because the two increments have stock units and $H$ has inverse-square-stock units.

In our example, the two increments are $(-2,0,2)$ and $(0,-2,2)$ litres. With $H$ equal to one quarter on each diagonal in inverse-square-litres, their cross term is one. That accounts exactly for $6=4+3-1$.

The correction need not always have this sign. Opposing increments can cancel and yield a negative cross term. For actions on disjoint coordinates in this diagonal model, the cross term is zero and singleton field values add. If a later potential couples coordinates, physical disjointness alone need not ensure disjoint potential support.

These are facts about the specified increments and potential. If a resolver changes action quantities depending on which companions are present, one must evaluate the resulting actual joint map. It would be wrong to reuse a fixed-increment proof after changing its inputs.

\section{What a path integral means}
An integral adds continuously varying small contributions. For our purpose, imagine a straight interpolation from the starting state to the joint endpoint:
\begin{equation*}
 y(\lambda)=z+\lambda\delta x_G,\qquad 0\leq\lambda\leq1.
\end{equation*}
The dimensionless parameter $\lambda$ marks a fraction of the whole change. Zero is the start, one is the endpoint. It is not automatically clock time.

The symbol $\nabla V$ means the gradient: the list of partial derivatives of $V$, equal to $\mu$ here. The rate of change of $V(y(\lambda))$ with respect to $\lambda$ is $\nabla V(y(\lambda))^{\mathsf T}\delta x_G$. Adding that rate along the complete path recovers the endpoint change. This is the fundamental theorem of calculus applied to the composed function. Reversing the sign gives
\begin{equation}\label{eq:path}
 E_G=-\int_0^1\nabla V(z+\lambda\delta x_G)^{\mathsf T}
                  \delta x_G\,d\lambda.
\end{equation}
It holds for a fixed continuously differentiable potential along the declared path. The integral and endpoint difference are two exact expressions for the same field value. Process burden, when present, must be accounted separately.\cite{foundation,gaussian}

For one transfer of quantity $q$, the same idea reads $E_e(q)=\int_0^q f_e(z+sS_e)\,ds$. The variable $s$ has stock units. Instead of multiplying all units by the initial field, the integral adds the changing field across the quantity interval. In the quadratic case it reproduces the finite correction from Chapter~\ref{ch:finite}.

\section{A common-path attribution}
Once the common straight path is declared, we can define a mathematical contribution for each fixed child increment:
\begin{align}
 R_a&=-\int_0^1\nabla V(z+\lambda\delta x_G)^{\mathsf T}
                 \delta x_a\,d\lambda\label{eq:attribution}\\
 &=-\mu(z)^{\mathsf T}\delta x_a
   -\frac12\delta x_a^{\mathsf T}H\delta x_G.\nonumber
\end{align}
The second line uses the quadratic gradient $\mu(z)+\lambda H\delta x_G$ and the fact that the integral of $\lambda$ from zero to one is one half. $R_a$ has the same account unit as $E_G$ and can be positive, zero or negative. The subscript labels an action, not a legal owner.

Add the first line over all children. The sum of their increments is $\delta x_G$, so linearity of multiplication and integration gives
\begin{equation}\label{eq:closure}
 \sum_{a\in G}R_a=E_G\qquad(C_G=0).
\end{equation}
This is exact mathematical closure under the stated assumptions. It is an Aumann--Shapley-type common-path marginal construction, a relationship already acknowledged in the historical foundation. It is not made unrelated to established allocation mathematics by giving its lines a new name.\cite{foundation,gaussian}

For the three-tank example, the initial marginals are $(1,0.5,-1.5)$ inverse litres. The initial linear contributions are five for $a$ and four for $b$. Each has a common-path curvature subtraction of $1.5$, giving $R_a=3.5$ and $R_b=2.5$. Their sum is six, the exact group total.

\section{Closure is not entitlement}
The straight line could describe a physical constant-rate action-time model, but only if that model is explicitly declared and physically valid. If only endpoints are specified, it is an interpolation convention for accounting. The arithmetic is the same; the interpretation is not.

Stock nonnegativity and upper capacities form a convex domain, so a straight segment between feasible endpoints remains within those particular bounds. That observation does not prove that real pipes, timing, intermediate stores or other action constraints realise the segment. A physically possible endpoint is not a complete physical history.

Nor does closure show that $3.5$ and $2.5$ are unique causal contributions, fair shares, legal property or permission to spend. An ownership map, an institutional allocation rule and a capacity policy are additional declarations. Identifying actual causal contributions requires the appropriate intervention or causal model, not merely an integral that sums correctly.

For comparison, if $a$ physically precedes $b$, its field value is four and the later value of $b$ is two. In the opposite order, $b$ gets three and the later $a$ gets three. Each sequence totals six because it reaches the same endpoint under fixed quantities. Different line allocations do not contradict the endpoint identity. They reveal different conventions or histories.

If real ordering changes feasibility, duration or the final state, the comparison must name those changes. ``Sequential'' cannot be an unnamed universal alternative. The historical sequential-parallel bridge makes that comparison discipline explicit.\cite{bridge}

\practice{At the reference, two ideal increments are exact opposites. What are the group's net increment and field value?}{The net increment is zero and the field value is zero. Each nonzero singleton would have negative value at the reference. Real resource use, wear or losses cannot be erased by cancellation of the displayed stock change; they would need representation in a more complete account.}

\carry{We can account for simultaneous groups without imposing a single-action world. Next we collect the exact conclusions, including sequences and cycles, and distinguish them from claims about recovery.}

\chapter{What follows mathematically, and what does not}\label{ch:limits}
\question{Which conclusions follow from definitions, and which need more evidence?}

A proof is only as broad as its assumptions. This chapter collects the exact conclusions and their limits. These are mathematical statements about a declared model; its physical adequacy and institutional use need separate justification.

\section{A small theorem ledger}
\textit{Affected-support cancellation.} Under one fixed additive potential, unchanged terms cancel. The local account must include every affected factor, including coupled factors (Chapter~\ref{ch:potential}).

\textit{Exact quadratic expansion.} Fixed references and positive scales give equation~\eqref{eq:finite}. Its positive correction for a nonzero increment makes the initial tangent overstate finite field reduction.

\textit{Endpoint/path equality.} For a fixed continuously differentiable potential on a declared differentiable path, integrating its rate of change gives its endpoint difference. Equation~\eqref{eq:path} assumes fixed additive increments; it does not establish a physical history.

\textit{Common-path closure.} Under those assumptions, the $R_a$ lines sum to the group field value. Closure does not select owners or establish fairness.

The scoped historical foundation and Gaussian reconciliation support these results.\cite{foundation,gaussian} We now derive a sequence identity for complete joint transitions.

\section{How the middle states cancel}
Suppose a first authoritative transition takes $x^0$ to $x^1$, then a second takes $x^1$ to $x^2$. Superscripts here label positions in a sequence, not powers. Use the same potential $V$ throughout and no unrecorded external changes between these transitions. With nonnegative process burdens $C_0$ and $C_1$, the two values are
\begin{align*}
 E_0&=V(x^0)-V(x^1)-C_0,\\
 E_1&=V(x^1)-V(x^2)-C_1.
\end{align*}
Adding cancels the two occurrences of $V(x^1)$, one negative and one positive. For $N$ consecutive transitions, the same cancellation gives
\begin{equation}\label{eq:telescoping}
 \sum_{n=0}^{N-1}E_n=V(x^0)-V(x^N)-\sum_{n=0}^{N-1}C_n.
\end{equation}
Every successor must be the next predecessor. Every transition is accounted once, under the same ruler and complete relevant state. A transition here may be the coherent joint action of a simultaneous group. This is a directly stated endpoint identity; it does not silently rewrite the one-action restrictions of the historical foundation's theorem.\cite{bridge,foundation}

In the two-tank illustration, a first three-litre delivery changes potential from four to $0.25$, giving $3.75$. A later one-litre delivery from that endpoint reaches $(10,10)$ and gives $0.25$. The sum is four. Quoting the later litre as if the first delivery had never happened would reuse the wrong baseline.

\section{Closed cycles and splitting}
If the complete endpoint returns to the complete start, $x^N=x^0$, the potential difference vanishes. Equation~\eqref{eq:telescoping} then gives
\begin{equation}\label{eq:cycle}
 \sum_n E_n=-\sum_n C_n\leq0.
\end{equation}
A genuine cost-free cycle has zero total field value. Positive included process burdens make the net negative. This is a conditional accounting conclusion about a closed undriven cycle with fixed $V$, complete state, faithful quantities and no duplication.

For a hand calculation, move three litres from $(14,6)$ to $(11,9)$, then move the same three litres back. The field values are $+3.75$ and $-3.75$. Their sum is zero. The calculation does not assert that physically moving real water back and forth has no resource cost; it is within the declared ideal lossless world.

Splitting a fixed endpoint change into consecutive pieces leaves the field total unchanged under these assumptions. Process burdens may differ if the actual split requires extra activations. A sequential field identity is not permission to ignore such burdens or to sum duplicated same-baseline singleton quotes.

Returning one visible stock is not necessarily a complete cycle. Fuel consumption, damage or an altered equipment state could remain elsewhere in a wider account. A cycle claim must restore the complete represented state relevant to it, not a convenient projection.

\section{External changes and a changed ruler}
Let external drive take $x^n$ to $z^n$ before a joint action takes $z^n$ to $x^{n+1}$. Define its signed potential change as
\begin{equation*}
 D_{{\rm ext},n}=V(z^n)-V(x^n).
\end{equation*}
Adding and subtracting $V(x^n)$ in each action value gives
\begin{equation}\label{eq:driven}
 \sum_n E_n=V(x^0)-V(x^N)
 +\sum_n D_{{\rm ext},n}-\sum_n C_n.
\end{equation}
The drive term has the same account unit as $V$. It can be positive or negative. A disturbance that brings the state closer to reference has a negative potential injection. The signed external record is not automatically an actor reward.\cite{foundation}

A different issue arises if the reference or scales change. Then adjacent terms may be evaluated under different versions of $V$ and no longer cancel. The model-change effect needs its own record. Rewriting earlier receipts under a new reference would hide that difference rather than resolve it.

These qualifications explain why the cycle identity is not a universal anti-gaming theorem. It does not settle fraud, omitted effects, coalition strategy, changed baselines or all driven histories. Historical adverse findings remain evidence about their original models. They are not erased by a cleaner exposition.

\section{A short prospective capacity bridge}
The Gaussian programme proposes persistent nonnegative experimental balances $B_i$ assigned to declared owners. These balances would constrain which actions a separately specified actor can afford. They are not physical water, installed pump capacity, money or a person's entitlement to care.

For the ideal $C=0$ proposal, suppose signed balance changes sum exactly to the joint action field reduction:
\begin{equation*}
 \sum_i\Delta B_{i,n}=V(z^n)-V(x^{n+1}).
\end{equation*}
Let an external audit $J$ accumulate the drive term, so $J_{n+1}=J_n+D_{{\rm ext},n}$. Add these declared updates to obtain
\begin{equation}\label{eq:capacity}
 V(x^{n+1})+\sum_i B_{i,n+1}-J_{n+1}
 =V(x^n)+\sum_i B_{i,n}-J_n.
\end{equation}
All terms use compatible dimensionless account units. The constant across events follows by cancellation. It depends on the stated update rules; it does not prove those rules desirable. Ownership, same-event netting and affordability are policy definitions beyond the identity. The current proposal introduces no borrowing, refill or freely transferable global pool.\cite{gaussian}

External drive changes the physical state and its audit; it does not itself earn actor capacity. A fixed audit after one disturbance still permits cycling, trapping or refusal. Under continuing drive, the physical state might have one kind of statistical behaviour while the balance and audit variables drift. A claim about the whole process would have to address them all.

\section{Why boundedness proves so little here}
In a closed nonnegative world with finite mass $M$, every stock satisfies $0\leq x_i\leq M$. Fixed finite references and positive scales therefore imply
\begin{equation}\label{eq:bound}
 V(x)\leq\frac12\sum_i
 \frac{\max\{(x_i^\star)^2,(M-x_i^\star)^2\}}{\sigma_i^2}<\infty.
\end{equation}
For each cell, the largest squared displacement on the interval occurs at one of its endpoints. Adding those bounds gives the displayed result, which may be loose because not all endpoint choices are jointly possible. It is still enough to establish boundedness.\cite{gaussian}

That bound holds before a beneficial EBU effect has been shown. A comparison policy in the same physical domain inherits it too. The interesting research questions concern recovery, time away from reference, repeated refusal, distribution and exposure, not unbounded growth of this fixed-domain potential.

\practice{A future study reports that both its EBU and comparison arms remained bounded. Has affordability been shown to cause stability?}{No. The finite-mass domain already supplies a bound. A stability, attraction or recovery claim requires the relevant definition, comparator, assumptions and evidence. An accounting invariant also does not supply the missing dynamics.}

\carry{The identities are useful because their meaning is precise. They bring us to the beginning of a research programme, not the end of one.}

\chapter{The equation as a beginning, not a promise}\label{ch:beginning}
\question{What can we now say precisely, and what remains open?}

Return to the clinic before its first patient arrives. The pump, the payment record and the need for safe care are still there. We have not replaced them with a formula. We have learned to ask what part of the event a formula can describe and what work is required before its number deserves an interpretation.

A physical action has a declared boundary and measurable quantities. A feasible action respects the constraints of the represented world. A reference and scale supply an explicit ruler. A potential measures the declared condition. A marginal describes a local sensitivity. The finite account compares complete endpoints and includes separately justified residual process burden once. A policy then uses or ignores that information under its own rules. Actors choose; execution and audit establish what occurred.

\section{The clinic revisited}
In the teaching world, delivering three litres from $(14,6)$ to $(11,9)$ reduced the declared potential by $3.75$. We can explain that number, reproduce it and show why a larger delivery can have a different sign. We can also explain why the unchanged surroundings cancel within an additive model.

We cannot infer from it that the clinic is safe, every patient is served or a real pump used no energy. Those would require more state variables, measurements and evidence. Nor does the positive number identify which institution should own a receipt, who pays the electricity bill or what a technician should earn.

Suppose the pump needs repair. The repair's immediate represented physical effects belong in its current account. If subsequent measurement establishes a changed capacity or efficiency, certification can support a new model version for future actions. Old records retain the version under which they were made. The problem of identifying and certifying persistent improvement remains a domain task; the equation does not pay immediately for an unsupported forecast.

An institution must separately authorise the work, carry any assigned settlement responsibility and compensate the worker. A negative current physical account is not personal blame. A beneficial future field change does not prove that the same worker will capture all its benefits. Public provision and access guarantees therefore remain substantive institutional questions.

Distance deserves the same discipline. A nearby service may require less transport, but proximity has no automatic EBU coefficient. A distant route may be more efficient under a particular physical account. The comparison must follow represented quantities, losses, time and capacity rather than an unexamined preference for either local or distant provision.

\section{The proposed experiment, briefly}
\textit{Prospective model proposal, not an executed result.} The Local Gaussian programme would specify external physical forcing; a frozen state; EBU-blind candidate actions or groups; joint physical feasibility; exact valuation and declared attributions; owner-specific capacity affordability; unweighted random actor choice; exact selected execution; and settlement and audit. No hidden value optimiser, service-first fallback or restoring dynamics supplies success in advance.\cite{gaussian}

A physical-random comparison would use declared physical rules without the capacity-affordability restriction. The future protocol must define action menus, owners, empty choices, numerical rules, disturbances and outcome measures before inspecting results. Sharing random draws would not automatically ensure equal applied disturbances after states diverge and feasibility differs.

The account might help, make little difference or perform poorly under that study. An introductory book cannot choose among those outcomes. It can make the question clear enough that the study is capable of answering it. No parameters, seeds, launch plan or experimental implementation are supplied by this edition.

\section{Where the other books belong}
The preserved Part II remains the home of rigorous mathematics under its named models. Its threshold potentials and controller proofs keep their original scope. A Gaussian specialisation is an additional declared construction, not a retroactive change to those proofs.

Part III preserves implementation and evidence history. D0, P1C and other historical findings must be read with their own designs and manifests. They are not Gaussian experimental evidence. This new volume changes the route into the subject without rewriting that record.

The current series architecture gives later volumes distinct responsibilities. Part IV concerns measurement, evidence and falsification. Part V concerns recovery and stochastic dynamics. Part VI develops multiple actions, interaction and local structure. Part VII follows routes, transformations and boundary effects across distance. Part VIII studies optional coordination, memory and feedback mechanisms. Part IX addresses institutions, rights, access, responsibility and transition.\cite{series}

Those responsibilities are research tasks, not guaranteed chapters of favourable results. In particular, a proof of accounting closure does not remove the need for causal identification, an acceptable ownership rule or a defensible treatment of people who cannot earn experimental capacity.

\section{What the reader can now challenge}
A careful reader can ask whether an action is physically defined, whether quantity and rate have been confused, whether the reference is compatible and whether it is reachable. They can check the sign convention, demand a finite endpoint, identify a missing affected factor and refuse to count the same burden twice.

They can also ask whether a plotted curve is an equation, an observation or a simulation. They can distinguish an unchanged total from a recovered function and a path allocation from a right to payment. These habits make the proposal more testable. They do not require faith in its eventual success.

The central question remains whether this explicit account is useful in worlds that can be measured and institutions that people can justify. The evidence of ecological pressures and unequal access makes that question worth investigating. The evidence needed to answer it is still separate.

We finish with the account whose meaning we can now state:
\begin{equation*}
 \boxed{\Delta E_a=V_A(z)-V_A(T_a(z))-C_a.}
\end{equation*}
One declared baseline. One coherent physical transformation. Complete affected support. A fixed, justified ruler. Compatible units and no double counting. The calculation is a beginning: EBU calculates, actors choose, and evidence must establish what follows.

\backmatter
\chapter*{Glossary}\addcontentsline{toc}{chapter}{Glossary}
\markboth{Glossary}{Glossary}
These definitions belong to the scope used in this book. A term can have a different technical meaning in another discipline; the model and context must travel with it.

\begin{description}
\item[Action.] A specified physical transformation, with its quantities, boundary and conditions. A proposed action is not evidence of an executed event.
\item[Actor.] The person, organisation or explicitly declared decision rule choosing an action. A value calculation does not itself make that choice.
\item[Affected support.] Every potential factor whose value can change in the action. It can extend beyond the immediately touched stocks in a coupled model.
\item[Attraction.] Approach towards a specified set or state under a specified dynamic law. A minimum of a function alone supplies no attraction.
\item[Attribution.] A declared assignment of a joint account to component actions. Mathematical closure does not establish causality, fairness or ownership.
\item[Baseline.] The state used for an action comparison. In a driven event, the frozen action baseline follows any declared prior external change.
\item[Boundedness.] Remaining within finite bounds. In this book's fixed-mass domain, stock and potential bounds already follow from the physical constraints.
\item[Capacity, experimental.] A prospective account balance limiting affordability under a separately declared rule. It is distinct from physical storage or throughput capacity.
\item[Conservation.] Constancy of a specified total within a declared boundary. A conserved total can be distributed in a functionally unsuitable way.
\item[Curvature.] How a slope changes as its input changes. In the quadratic potential it supplies an exact finite correction to the initial marginal estimate.
\item[Derivative.] The limiting rate of change of a function with respect to one input. Its unit is the function's unit divided by the input's unit.
\item[Dissipation.] Irreversible physical degradation of available driving differences. Its representation and valuation require a declared physical model; it is not a synonym for every monetary expense.
\item[Drive.] A declared external change to the represented physical state. Its potential effect must be distinguished from the actor's subsequent action value.
\item[EBU.] The project's name for a proposed account of physical action. Its core finite expression is an affected potential difference minus a separately justified process burden. The letters do not establish an energy unit or a currency.
\item[Equilibrium.] A state left unchanged by a specified dynamic law. Thermodynamic equilibrium has the more specific physical meaning discussed in Chapter~\ref{ch:balance}.
\item[Externality.] An effect on others that is not fully reflected in the private costs or benefits guiding a decision.
\item[Feasibility.] Satisfaction of the declared physical constraints. It does not by itself establish permission, affordability or desirability.
\item[Field, directional.] The initial marginal reduction of potential per unit quantity in a specified action direction. It can change over a finite action.
\item[Functional balance.] Conditions compatible with specified functioning, potentially maintained through continuous throughput. It need not mean equal stocks or immobility.
\item[Gaussian ruler.] Here, a quadratic measure of deviation associated with the shape of a Gaussian reference density. It does not assert normally distributed physical trajectories.
\item[Gradient.] The vector of partial derivatives of a scalar function. In this separable model it is the vector of local marginal burdens.
\item[Hessian.] The matrix of second derivatives. For the fixed Gaussian ruler it is diagonal, with positive entries $1/\sigma_i^2$.
\item[Homeostasis.] Active regulation of selected internal variables within ranges compatible with function. It depends on mechanisms and resources.
\item[Identity.] An equality following from stated definitions and assumptions. Its exactness does not demonstrate the empirical adequacy of those assumptions.
\item[Local information.] Information available within a declared physical or informational neighbourhood. Complete affected support and limited actor knowledge are different requirements.
\item[Lossless transfer.] A transfer in which source decrease equals destination increase for the represented resource. Real losses require explicit physical accounting.
\item[Marginal burden.] The sensitivity of potential to one stock, holding other coordinates fixed. Its sign indicates the local slope, not a moral judgement.
\item[Normal distribution.] A probability distribution with a bell-shaped density. Choosing its geometry as a ruler is distinct from fitting it to observations.
\item[Permission.] An institutional or account rule governing which feasible actions may proceed. It must be stated separately from valuation.
\item[Potential.] The declared scalar function of physical state used for valuation. In this book it is dimensionless deviation, not automatically thermodynamic energy or social welfare.
\item[Process burden.] A separately declared nonnegative physical contribution in compatible account units, subject to the no-double-counting boundary. The ideal lossless illustration sets it to zero.
\item[Public good.] In the ideal economic definition, a benefit that is nonrival and nonexcludable. Actual services may have congestion and exclusion constraints.
\item[Recovery.] A return of specified variables or functions to a declared condition after disturbance, with a stated criterion and time horizon. It requires dynamics and evidence beyond an account identity.
\item[Reference.] The declared comparison state. Compatibility with mass and capacities does not establish reachability or equilibrium.
\item[Scale.] A fixed positive stock quantity setting the ruler's sensitivity to displacement. It is not automatically an observed standard deviation or a safety limit.
\item[Settlement.] Application of separately declared rules to account balances after an event. Field value, compensation and personal responsibility are distinct questions.
\item[Stability.] A specified form of controlled response to perturbation under a specified dynamic law. One must name the relevant definition; boundedness alone is insufficient.
\item[Steady state.] Constancy of represented stocks or macroscopic variables over time. It may coexist with continuing input and output flows.
\item[Telescoping.] Cancellation of intermediate endpoint terms in a sequence of differences evaluated under one fixed potential.
\item[Throughput.] Physical flow through a system over time. Constant stock does not imply zero throughput.
\end{description}

\chapter*{Symbols and units}\addcontentsline{toc}{chapter}{Symbols and units}
\markboth{Symbols and units}{Symbols and units}
The stock unit in the illustrations is the litre (L). A different homogeneous resource may use another stock unit. Account quantities below are dimensionless for the Gaussian construction. Superscript $\star$ labels a reference; $\mathsf T$ transposes a vector; superscript $n$ labels an event state. None of these means a square.

\renewcommand{\arraystretch}{1.16}
\begin{longtable}{@{}p{65pt}p{194pt}p{80pt}@{}}
\toprule Symbol & Meaning & Unit here\\\midrule\endfirsthead
\toprule Symbol & Meaning & Unit here\\\midrule\endhead
\bottomrule\endfoot
$i,j$ & Cell indices & Labels\\
$x_i,\ x$ & Stock at a cell; full stock vector & L\\
$z$ & Frozen post-drive action baseline & L\\
$x_i^\star$ & Declared reference stock & L\\
$\sigma_i>0$ & Fixed scale & L\\
$M$ & Conserved total stock & L\\
$K_i$ & Physical storage capacity & L\\
$a,\ e,\ G$ & Action, directed edge, action group & Labels/set\\
$q$ & Finite transferred quantity & L\\
$r,\ \Delta t$ & Transfer rate; duration & L/min; min\\
$e_i$ & Unit vector selecting cell $i$ & Dimensionless\\
$S_e=-e_i+e_j$ & Lossless transfer direction & Dimensionless\\
$\delta x_a$ & Physical increment of action $a$ & L\\
$\delta x_G$ & Sum of declared group increments & L\\
$T_a(z)$ & State after action map $a$ & L\\
$A$ & Complete affected potential support & Set\\
$V_i,\ V,\ V_A$ & Local, total and affected potential & Account unit\\
$\mu_i$ & $\partial V/\partial x_i$: marginal burden & L$^{-1}$\\
$\mu=\nabla V$ & Vector of partial derivatives & L$^{-1}$\\
$f_e=\mu_i-\mu_j$ & Directional marginal field & L$^{-1}$\\
$H$ & Fixed Hessian, diagonal $1/\sigma_i^2$ & L$^{-2}$\\
$p_i(u)$ & Gaussian reference density & L$^{-1}$\\
$C_a$ & Separate residual process burden & Account unit\\
$E,\ \Delta E_a$ & Finite signed action account & Account unit\\
$E_G$ & Joint field reduction (zero cost) & Account unit\\
$E_a^{(0)}$ & Singleton value at common start & Account unit\\
$R_a$ & Declared common-path attribution & Account unit\\
$\lambda$ & Fraction along straight interpolation & Dimensionless\\
$n,\ N$ & Cell count (Ch.12); event index and transition count (Ch.17) & Counts\\
$D_{{\rm ext},n}$ & Signed external potential injection & Account unit\\
$B_i$ & Prospective experimental balance & Account unit\\
$J$ & Signed cumulative external audit & Account unit\\
\end{longtable}
\renewcommand{\arraystretch}{1}

In $u^{\mathsf T}v$, multiply matching vector entries and add them. In $\sum_i$, add over the declared indices. In $\int$, accumulate the specified continuously varying quantity over its stated interval. These operations are introduced where first needed; they do not select an actor policy.

\chapter*{Sources and further reading}\addcontentsline{toc}{chapter}{Sources and further reading}
\markboth{Sources}{Sources}
The factual references below support the selected claims, not EBU's effectiveness. Dates and scopes matter: the climate figures are from a 2023 assessment, and their observation and projection periods differ. Web sources were checked on 19 September 2026. The accompanying claim/source ledger records exact passages, retrieval limits and unused claims. Numbered S references concern external evidence or definitions; M references identify the project's scoped mathematical and editorial sources.

\begingroup
\renewcommand{\chapter}[2]{}
\begin{thebibliography}{M04}
\markboth{Sources}{Sources}
\raggedright
\bibitem[S01]{boe} McLeay, M., Radia, A. and Thomas, R. (14 March 2014). \textit{Money creation in the modern economy}. Bank of England Quarterly Bulletin, 54(1), 14--27. Overview and lending constraints. \href{https://www.bankofengland.co.uk/quarterly-bulletin/2014/q1/money-creation-in-the-modern-economy}{Bank of England publication and PDF}.
\bibitem[S02]{externalities} Helbling, T. \textit{Externalities: Prices Do Not Capture All Costs}. IMF, Finance \& Development, Back to Basics; maintained explainer, no date assumed. Definitions and public-goods discussion. \href{https://www.imf.org/en/publications/fandd/issues/series/back-to-basics/externalities}{IMF article}.
\bibitem[S03]{ecb} European Central Bank (2025). \textit{Monetary policy strategy statement}. June 2025 version, sections 2 and 8--10, pp.1--2. \href{https://www.ecb.europa.eu/mopo/strategy/strategy-review/ecb.strategyreview202506_strategy_statement.en.pdf}{Official strategy statement}.
\bibitem[S04]{ets} European Commission. \textit{About the EU ETS}. Maintained institutional page; ``How does the EU ETS work?'' No current price or effectiveness estimate used. \href{https://climate.ec.europa.eu/areas-action/carbon-markets/about-eu-ets_en}{Commission overview}.
\bibitem[S05]{water} Food and Agriculture Organization. \textit{Water scarcity}. Maintained Land \& Water overview; three-type classification. \href{https://www.fao.org/land-water/water/water-scarcity/en}{FAO overview}.
\bibitem[S06]{health} World Health Organization (6 May 2025). \textit{World report on social determinants of health equity}. ISBN 978-92-4-010758-8. Qualitative claim based on the official overview, not an independent report replication. \href{https://www.who.int/publications/i/item/9789240107588}{WHO publication overview}.
\bibitem[S07]{mto} Chetty, R., Hendren, N. and Katz, L.F. (2016). \textit{The Effects of Exposure to Better Neighborhoods on Children: New Evidence from the Moving to Opportunity Experiment}. American Economic Review, 106(4), 855--902. \href{https://doi.org/10.1257/aer.20150572}{doi:10.1257/aer.20150572}. The bounded description in Chapter~\ref{ch:access} follows the \href{https://opportunityinsights.org/paper/newmto/}{authors' study summary} (earlier working-paper summary, May 2015). US programme recruitment, 1994--1998; no universal effect size inferred.
\bibitem[S08--09]{ipcc} IPCC (2023). \textit{Climate Change 2023: Synthesis Report. Summary for Policymakers}. In H. Lee and J. Romero (eds.), AR6 Synthesis Report. Observed warming: A.1--A.1.1 and note 5, printed p.4. Conditional 2081--2100 warming: B.1.1, notes 28 and 30, p.12. Increasing risks: B.2, p.14. \href{https://www.ipcc.ch/report/ar6/syr/downloads/report/IPCC_AR6_SYR_SPM.pdf}{Official SPM PDF}.
\bibitem[S10]{ipbes} IPBES (2019). \textit{Summary for policymakers of the global assessment report on biodiversity and ecosystem services}. S. D\'iaz et al. (eds.). A5 (printed pp.11--12) and background paragraph 6 (p.24) supply the threatened-species estimate and its uncertainty. \href{https://files.ipbes.net/ipbes-web-prod-public-files/2020-02/ipbes_global_assessment_report_summary_for_policymakers_en.pdf}{Official published SPM PDF}. This is the 2019 assessment, not a current species census.
\bibitem[S11]{resources} UNEP and International Resource Panel (1 March 2024). \textit{Global Resources Outlook 2024}. The conditional 2020--2060 extraction comparison is taken from the \href{https://www.unep.org/resources/Global-Resource-Outlook-2024}{official report overview}; the overview supplies no formal interval for that headline.
\bibitem[S12]{air} World Health Organization (24 October 2024). \textit{Ambient (outdoor) air pollution}. Fact sheet; burden estimates refer to 2019. Only the qualitative health-burden finding is used. \href{https://www.who.int/news-room/fact-sheets/detail/ambient-(outdoor)-air-quality-and-health}{WHO fact sheet}.
\bibitem[S13]{seea} United Nations Statistics Division. \textit{Ecosystem Accounting} and \textit{SEEA Implementation Guide: Introduction}. Adoption in March 2021; physical accounts and the distinction between statistical standards and monetary-valuation recommendations. \href{https://seea.un.org/en/methodology/ecosystem-accounting}{Framework overview}; \href{https://seea.un.org/en/content/seea-implementation-guide-introduction}{implementation guide}. Official indexed passages were inspected where direct retrieval was restricted.
\bibitem[S14]{homeostasis} McFarland, J. et al. (June 2016). \textit{A conceptual framework for homeostasis: development and validation}. Advances in Physiology Education, 40(2), 213--222, especially framework statements H1.4--H1.8. \href{https://doi.org/10.1152/advan.00103.2015}{doi:10.1152/advan.00103.2015}; \href{https://pmc.ncbi.nlm.nih.gov/articles/PMC5002438/}{open article}. A physiological teaching framework, not an economic model.
\bibitem[S15]{adaptation} Lan, G., Sartori, P., Neumann, S., Sourjik, V. and Tu, Y. (25 March 2012). \textit{The energy-speed-accuracy trade-off in sensory adaptation}. Nature Physics, 8, 422--428. \href{https://doi.org/10.1038/nphys2276}{doi:10.1038/nphys2276}. The use here is limited to the publisher abstract's theory and bacterial-adaptation scope.
\bibitem[S16]{normal} NIST/SEMATECH. \textit{e-Handbook of Statistical Methods}, section 1.3.6.6.1, Normal Distribution. Maintained mathematical reference, no publication date assumed. \href{https://www.itl.nist.gov/div898/handbook/eda/section3/eda3661.htm}{Density definition and assumptions}.
\bibitem[M01]{foundation} Grzyb, K., EBU project. \textit{V3.0 Local EBU Foundation Draft}, repository source \nolinkurl{V3.0_LOCAL_EBU_FOUNDATION_DRAFT.md}. Sections 6--8 and selected common-path construction in section 10, with the restrictions recorded in the claim ledger. Last source-changing commit: \texttt{d45f059b2471}. Historical model assumptions remain in force.
\bibitem[M02]{gaussian} Grzyb, K., EBU project. \textit{Local Gaussian EBU Programme Reconciliation}, repository source \nolinkurl{LOCAL_GAUSSIAN_EBU_PROGRAMME_RECONCILIATION.md}, especially sections 3--7 and 10. Prospective definitions and research questions, not executed experimental findings. Version locked at repository HEAD \texttt{05f3cb02d0c4}.
\bibitem[M03]{bridge} Grzyb, K., EBU project. \textit{Sequential--Parallel Bridge}, repository source \nolinkurl{SEQUENTIAL_PARALLEL_BRIDGE.md}, sections 3.8 and 4--7. Last source-changing commit: \texttt{2676912a3d16}. Comparisons must retain their named maps, quantities, ordering and scope.
\bibitem[M04]{series} Grzyb, K., EBU project. \textit{EBU Future Books Structure} and \textit{Local Gaussian EBU Book-Series Reconciliation}. Repository sources \nolinkurl{EBU_FUTURE_BOOKS_STRUCTURE.md} and \nolinkurl{LOCAL_GAUSSIAN_EBU_BOOK_SERIES_RECONCILIATION.md}; locked at HEAD \texttt{05f3cb02d0c4}. Later-volume responsibilities are prospective.
\end{thebibliography}
\markboth{Sources}{Sources}
\endgroup

\section*{The preserved baseline and editorial record}
The preserved \textit{EBP Book, Part I, Unified Explanatory Edition} has 296 PDF pages. It is a content baseline, not the editable source of this edition. Its SHA-256 begins \texttt{335ed5c6d3541d48}; the complete digest, identities of Parts II and III, and documentary locks are in the accompanying source manifest.

The new manuscript follows \textit{Book I Introduction Blueprint}, \textit{Book I Generation Handover} and the completed reconciliation at the locked repository version. A section-level disposition records retained ideas, rewritten explanations, condensed material, later-volume responsibilities and omissions from this introduction. Nothing was inserted into the preserved later volumes.

The tank quantities, diagrams and arithmetic were created for this book. They are labelled illustrations and have no empirical outcome status. The mathematical review notes identify the derivations and fixed-state arithmetic checked during preparation. No new model trajectory was produced.

\end{document}
